\documentclass[11pt]{amsart}
\usepackage[a4paper,margin=26mm]{geometry}
\usepackage{amsmath,amssymb,amsthm,mathtools,booktabs,array,longtable}
\usepackage[expansion=false]{microtype}
\usepackage[hidelinks,breaklinks]{hyperref}
\usepackage{xurl}
\newtheorem{theorem}{Theorem}[section]
\newtheorem{proposition}[theorem]{Proposition}
\newtheorem{lemma}[theorem]{Lemma}
\newtheorem{corollary}[theorem]{Corollary}
\theoremstyle{definition}
\newtheorem{definition}[theorem]{Definition}
\theoremstyle{remark}
\newtheorem{remark}[theorem]{Remark}
\newcommand{\Q}{\mathbb Q}
\newcommand{\Z}{\mathbb Z}
\newcommand{\C}{\mathbb C}
\newcommand{\F}{\mathbb F}
\newcommand{\HH}{\mathfrak H}
\newcommand{\nm}{\mathrm}
\DeclareMathOperator{\Sym}{Sym}
\DeclareMathOperator{\Ind}{Ind}
\DeclareMathOperator{\tr}{tr}
\DeclareMathOperator{\Aut}{Aut}
\DeclareMathOperator{\GL}{GL}
\DeclareMathOperator{\SL}{SL}
\newcommand{\workbenchurl}{https://github.com/KokunoYumeto/zeta-and-l-functions-workbench}
\newcommand{\recordurl}{https://github.com/KokunoYumeto/mathematics-workbenches}
\title[Moonshine after Alp\"oge's theorem]{Moonshine after Alp\"oge's theorem\\
\large Dirichlet shadows, off-line Epstein zeros, and CM reductions}
\author{}
\date{24 November 2026}
\begin{document}
\begin{abstract}
Alp\"oge's theorem establishes the Riemann hypothesis for Dirichlet
$L$-functions, bringing a theorem on zeros to the degree-one arithmetic
inside the moonshine of Conway--Norton, Mason, and Cheng--Duncan--Harvey.
We identify every unary-shadow character constituent with a rescaled
Dirichlet $L$-function. In contrast, Saias and Weingartner's theorem implies
infinitely many off-line zeros for every shadow mixing distinct primitive
characters, including the first component at lambency six. A theorem of
Nakamura and Pa\'nkowski also gives infinitely many off-line zeros of
every rank-twenty-four even unimodular Epstein zeta function, including
that of the Leech lattice.
We give exact theta coordinates for 129 lattice representatives, classify
the weight-one Conway eta forms, and prove the associated CM identities,
explicit class-number bounds, Fricke constants, and a degree-two isogeny
for the cusp cover of $X_0(6)$. The positive-infinity-type Hecke Riemann
hypothesis remains a separate question: a local cubic symmetry is available,
but the required global prime-ideal estimate is not established here.
The mathematics and text were produced by GPT-6 Astra (ChatGPT Pro mode);
the checking records distinguish computations and re-proofs by separate
instances of the same model, and no human has checked the paper.
\end{abstract}
\maketitle

\section{Introduction}
Alp\"oge's theorem of 20 November 2026 proves the Riemann hypothesis for
every Dirichlet $L$-function: zeros in the open critical strip have
real part $1/2$, and there are no zeros strictly to the right of that
line, with the principal-character pole convention excluded
\cite{AlpogeLean}.
We use that theorem as a published input. The independent checks reported
by the Lean community concern the statement, compilation, axioms, and
kernel replay of the formal proof \cite{LeanVerification}.
An earlier AI-produced result relevant to the extension problem is
OpenAI's zero-free half-plane $\Re s>7/8$ for finite-order Hecke
characters over $\Q(\sqrt{-3})$ and Dirichlet characters
\cite[Theorem~1.1, p.~4]{OpenAIQRH}. Its cubic-theta and completed-index
arguments concern the same arithmetic field. Positive infinity type
requires an archimedean character beyond the theorem's finite-order
hypothesis; the companion note compares this issue with the subsequent
unitary-character extensions.

The objects to which we apply it come from several distinct theorems.
Conway and Norton identified the modular functions conjecturally attached
to Monster conjugacy classes \cite[Sect.~2 and Table~2]{CN1979}; Borcherds
proved the corresponding moonshine theorem \cite[Theorem~1.1]{Borcherds1992}.
Mason attached Hecke eigenforms to the cycle shapes of $M_{24}$
\cite{Mason1985}. Cheng, Duncan and Harvey associated mock modular forms
and unary-theta shadows to the Niemeier root systems \cite{CDH2014}, and
Duncan, Griffin and Ono proved the existence of the umbral modules
\cite[Theorem~1.1]{DGO2015}. These results identify the forms; none asserts
a common Riemann hypothesis for all the scalar series made from them.

Here are the main analytic statements. Our normalization is
$q=e^{2\pi i\tau}$, and
$S_{m,r}=\sum_{n\equiv r\ (2m)}nq^{n^2/(4m)}$.
For $H=\sum_{n\geq1}c(n)nq^{n^2/(4m)}$, write
$L(H,s)=\sum_{n\geq1}c(n)n(n^2/(4m))^{-s}$.

\begin{theorem}[The degree-one part of every shadow]\label{thm:intro-dirichlet}
Let $m\geq2$, $d\mid2m$, and let $\chi$ be an odd Dirichlet character
modulo $2m/d$. For
$G_{d,\chi}=\sum_{r\bmod(2m/d)}\chi(r)S_{m,dr}$ one has
\[
 L(G_{d,\chi},s)=2d(4m/d^2)^sL(2s-1,\chi).
\]
Every fixed component of every umbral shadow, including every twining,
is a finite linear combination of these functions. Every zero of a
nonzero $G_{d,\chi}$ in $1/2<\Re s<1$ lies on $\Re s=3/4$.
\end{theorem}

The last assertion follows from Alp\"oge's theorem. The spanning
assertion does not transfer that conclusion to an arbitrary sum.

\begin{theorem}[Mixed shadows have off-line zeros]\label{thm:intro-mixed}
Let $H\ne0$ be a fixed linear combination of the $S_{m,r}$, and express
its periodic coefficient series uniquely as
\[
 \sum_{n\geq1}c(n)n^{-z}=\sum_{\psi\ {\nm primitive}}P_\psi(z)L(z,\psi),
\]
where the $P_\psi$ are Dirichlet polynomials. If at least two are nonzero,
there exists $\epsilon>0$ such that, for every
$3/4<\alpha<\beta\leq1+\epsilon$, the number of distinct zeros of
$L(H,s)$ in $\alpha<\Re s<\beta$, $|\Im s|\leq T$, is bounded above
and below by positive multiples of $T$ for sufficiently large $T$.
This applies to $5S_{6,1}+S_{6,5}$ and $S_{6,1}+5S_{6,5}$.
In particular, both have infinitely many zeros with $\Re s>1$.
\end{theorem}

This applies the Primitive Dirichlet Character Basis theorem
and the zero theorem of Saias and Weingartner \cite[pp.~335--336]{SW2009}.
It settles the single-component question whenever at least two primitive
characters occur. If only one occurs, the remaining Dirichlet polynomial
must be examined separately. An identically zero twining is excluded from
statements about zero sets.

\begin{theorem}[Off-line zeros for every Niemeier Epstein function]
\label{thm:intro-niemeier}
Let $L$ be a positive-definite even unimodular lattice of rank $24$,
and define
\[
 Z_L(w)=\sum_{v\in L\setminus\{0\}}
             \bigl((v,v)/2\bigr)^{-w}.
\]
For every $23/2<\alpha<\beta<12$, there are positive constants
$c_1,c_2,T_0$, depending on $L,\alpha,\beta$, such that for $T\geq T_0$
the number of zeros in
$\alpha<\Re w<\beta$, $|\Im w|\leq T$, counted with multiplicity,
lies between $c_1T$ and $c_2T$. The lower bound also holds for
distinct simple zeros. The reflected assertion holds in
$12-\beta<\Re w<12-\alpha$.
\end{theorem}

This is an application of the hybrid-universality theorem of
Nakamura and Pa\'nkowski \cite[Main Theorems~1--2 and Theorem~E,
pp.~4--5 of the preprint]{NP2016}. It is unconditional.
In particular, the scalar Epstein zeta function of the Leech lattice
does not satisfy a Riemann hypothesis on its functional-equation line
$\Re w=6$. The theorem locates infinitely many zeros; it does not
classify every zero or decide whether any lies in $\Re w>12$.

The finite-dimensional structures are explicit. At lambency six the
twining characters form symmetric powers of a faithful two-dimensional
representation of $\GL_2(\F_3)$. The cyclic residue action on the rational
space of the Niemeier lattice with roots $A_5^4D_4$ gives
$\Delta$, $\eta(2\tau)^{12}$, $\eta(3\tau)^8$, and $\eta(6\tau)^4$.
Their CM members lead to Hecke characters of $\Q(\sqrt{-3})$ of positive
infinity type. These are not Dirichlet characters. The Hecke Riemann
hypotheses for those forms are not proved here.

The motivating lattice formulas and residue action were developed in
the study of the $(3,4,\infty)$ family of complex six-spheres and its
arithmetic lattice neighbours; a subsequent moonshine atlas brought
those formulas together. We credit the relevant statements
in their respective sections. Their earlier numbering and the complete
check records are indexed at \url{\recordurl}. The accompanying programs
are collected at the workbench \url{\workbenchurl}.
No claim of priority is made.

\section{Proof overview}
The analytic distinction is elementary before the theorem on zeros is
used. A unary shadow has a periodic coefficient function. Dividing its
indices according to their greatest common divisor with the modulus
reduces that function to functions on finite unit groups. Character
orthogonality gives Dirichlet $L$-functions, with every dilation and
omitted Euler factor visible. The primitive-character expansion is unique;
that is what allows Saias and Weingartner's theorem to decide whether
mixing characters creates zeros even in the half-plane of absolute
convergence.

The rank-twenty-four Epstein problem has a different mechanism.
Its theta series is a sum of one Eisenstein series and a multiple of
Ramanujan's cusp form. After shifting the variable by eleven, the
cusp-form term converges absolutely throughout the universality strip
of the remaining zeta factor. Hybrid universality approximates a
nonvanishing target chosen so that the sum has a simple zero.
Rouch\'e's theorem then turns those approximations into a positive
linear density of simple zeros. The functional equation reflects
them to the opposite edge of the critical strip.

The lattice and CM arguments use different information. A Coxeter
eigenvalue computation gives the Frame shapes. Complex multiplication
identifies two eta products and decomposes the symmetric cube;
local inertia determines the factors at the bad primes. For theta series,
small modular-form spaces turn exact shell counts into identities.
Fricke constants follow from the transformation of $\eta$ and the
determinant of a root lattice. A singular-modulus test decides the CM
zeros of Niemeier theta series.

The class-number estimates require signs and normalization to be kept.
The class-number contribution to a virtual multiplicity can be negative
or zero, and Tunnell's proportionality constant varies with the
discriminant. The positive-infinity-type Hecke problem requires a global
estimate on prime ideals. A geometric symmetry proves the necessary local
commutation, but does not supply that estimate.

% New exposition and proofs. No source-paper text is reproduced.
\section{Dirichlet constituents and zeros of umbral shadows}
\label{sec:shadow-dirichlet}
This section identifies exactly which scalar functions obtained from an umbral
shadow inherit the Dirichlet Riemann hypothesis.  It also determines the zeros
of every component involving two distinct primitive Dirichlet characters by
applying the theorem of Saias and Weingartner \cite[pp.~335--336]{SW2009}.
In particular, the first component for $A_5^4D_4$ has infinitely many zeros to
the right of its expected critical line, and even in its half-plane of absolute
convergence.  This last conclusion is unconditional.

Put $q=e^{2\pi i\tau}$ and, for $m\geq2$ and $r\in\mathbb Z/2m\mathbb Z$, put
\begin{equation}\label{eq:unary-definition}
 S_{m,r}(\tau)=\sum_{\substack{n\in\mathbb Z\\n\equiv r\pmod{2m}}}
 n q^{n^2/(4m)}.
\end{equation}
Thus $S_{m,-r}=-S_{m,r}$ and $S_{m,0}=S_{m,m}=0$.
This is the normalization of Cheng, Duncan and Harvey
\cite[Sect.~3.4, equation (3.39), p.~37]{CDH2014}; in particular, its exponent denominator
is $4m$.  For a series
\[
 H(\tau)=\sum_{n\geq1}c(n)nq^{n^2/(4m)}
\]
with periodic coefficients we use the Dirichlet transform
\begin{equation}\label{eq:shadow-L-normalization}
 L(H,s)=\sum_{n\geq1}c(n)n\left(\frac{n^2}{4m}\right)^{-s}
       =(4m)^sD_c(2s-1),\qquad
 D_c(z)=\sum_{n\geq1}\frac{c(n)}{n^z}.
\end{equation}
The defining sums converge absolutely for $\Re s>1$ and $\Re z>1$,
respectively.  They thereafter mean their meromorphic continuations.

\begin{theorem}[Character constituents at every lambency]\label{thm:all-shadow-characters}
Let $N=2m$, let $d\mid N$, put $Q=N/d$, and let $\chi$ be an odd Dirichlet
character modulo $Q$.  Define
\[
 G_{d,\chi}(\tau)=\sum_{r\bmod Q}\chi(r)S_{m,dr}(\tau).
\]
Then
\begin{equation}\label{eq:all-shadow-character-L}
 L(G_{d,\chi},s)=2d\left(\frac{4m}{d^2}\right)^sL(2s-1,\chi).
\end{equation}
If $\chi^*$ is the primitive character inducing $\chi$, of conductor $f$, then
\begin{equation}\label{eq:imprimitive-shadow}
 L(2s-1,\chi)=L(2s-1,\chi^*)
 \prod_{\substack{p\mid Q\\p\nmid f}}
       \left(1-\chi^*(p)p^{1-2s}\right).
\end{equation}
Consequently the Dirichlet Riemann hypothesis implies that every zero of
this nonzero constituent in $1/2<\Re s<1$ satisfies $\Re s=3/4$.
The omitted Euler factors have their zeros on $\Re s=1/2$.
Every component of every fixed umbral shadow or twined shadow is a finite
linear combination of the $G_{d,\chi}$.
\end{theorem}
\begin{proof}
In the sum defining $G_{d,\chi}$, set $n=dk$.  Summing over residue classes
first gives
\[d\sum_{k\in\mathbb Z}\chi(k)kq^{d^2k^2/(4m)}.\]
Since $\chi(-1)=-1$, the positive and negative terms agree.  Mellin expansion
in the absolutely convergent half-plane proves
\eqref{eq:all-shadow-character-L}.  The usual Euler product proves
\eqref{eq:imprimitive-shadow} there and hence everywhere by continuation.
For a prime in that product $|\chi^*(p)|=1$, so a zero of its factor has
$p^{1-2\Re s}=1$.  The claimed zero lines follow.

For the spanning assertion, write a component as
$H=\sum_{r\bmod N}h_rS_{m,r}$.  Its coefficient in
\eqref{eq:shadow-L-normalization} is the odd periodic function
$c(n)=h_n-h_{-n}$.  On the stratum $\gcd(n,N)=d$, the function $k\mapsto c(dk)$
is an odd function on $(\mathbb Z/Q\mathbb Z)^\times$.  Orthogonality of
characters gives
\begin{equation}\label{eq:character-stratum-coefficients}
 c(dk)=\sum_{\chi\bmod Q}C_{d,\chi}\chi(k),\qquad
 C_{d,\chi}=\frac1{\varphi(Q)}
    \sum_{u\in(\mathbb Z/Q\mathbb Z)^\times}c(du)\overline{\chi(u)}.
\end{equation}
Changing $u$ to $-u$ shows that $C_{d,\chi}=0$ for even $\chi$.
Thus $H=\frac12\sum_{d,\chi}C_{d,\chi}G_{d,\chi}$.
The umbral and twined shadows are such linear combinations by their defining
matrix formula $S_g^X=\Omega_g^XS_m$
\cite[Sect.~4.1 and Sect.~5.1]{CDH2014}.  No restriction on the Coxeter number or on
$g$ was used.
\end{proof}

The distinction between a constituent and a sum of constituents is essential.
In particular, a primitive character may occur on several gcd strata, and
adding those strata produces a Dirichlet polynomial which need not be a
product of omitted Euler factors.

\begin{theorem}[Classification of zeros of a scalar shadow]\label{thm:shadow-zero-classification}
Let $H$ be a nonzero fixed linear combination of the $S_{m,r}$, and write its
coefficient function $c$ as above.  There is a unique finite expansion
\begin{equation}\label{eq:primitive-shadow-expansion}
 D_c(z)=\sum_{\psi\ {
m primitive}}P_\psi(z)L(z,\psi)
\end{equation}
in which each $P_\psi$ is a Dirichlet polynomial.  All characters occurring
are odd.

If at least two $P_\psi$ are nonzero, there is an $\epsilon>0$ such that,
for every $3/4<\alpha<\beta\leq1+\epsilon$, the number of distinct zeros of
$L(H,s)$ in
\[
 \alpha<\Re s<\beta,\qquad |\Im s|\leq T
\]
is bounded above and below by positive constant multiples of $T$, for all
sufficiently large $T$.  In particular there are infinitely many zeros in
$\Re s>1$.  These conclusions need no Riemann hypothesis.

If exactly one $P_\psi$ is nonzero, the zeros are those of
$P_\psi(2s-1)$ and $L(2s-1,\psi)$, with multiplicity, apart from no further
factors.  The Dirichlet Riemann hypothesis controls the second factor only.
\end{theorem}
\begin{proof}
Split the sum for $D_c$ by $\gcd(n,N)=d$ and use
\eqref{eq:character-stratum-coefficients}.  This yields
\[
 D_c(z)=\sum_{d\mid N}\sum_{\chi\bmod N/d}
 C_{d,\chi}d^{-z}L(z,\chi).
\]
Replacing each imprimitive $L$-function by its primitive one times its finite
Euler product gives \eqref{eq:primitive-shadow-expansion}.
Uniqueness is precisely the Primitive Dirichlet Character Basis theorem of
Saias and Weingartner \cite[p.~335]{SW2009}.  Thus having two nonzero
primitive terms is equivalent to $D_c$ not being a Dirichlet polynomial times
one Dirichlet $L$-function.  Their main theorem
\cite[pp.~335--336]{SW2009} supplies $\eta>0$ and upper and lower linear
zero-counting bounds in every strip
$1/2<\sigma_1<\sigma_2\leq1+\eta$ for $D_c$.
Set $\sigma_1=2\alpha-1$, $\sigma_2=2\beta-1$ and
$\epsilon=\eta/2$.  The affine change $z=2s-1$ doubles the height and
preserves multiplicity; the factor $(4m)^s$ has no zero.  This proves the
assertion.  In the one-character case $\psi$ is odd and hence nonprincipal,
so its $L$-function is entire.  The stated product therefore has no pole
which could cancel a zero.
\end{proof}

This theorem answers the question for every lambency at once whenever a
component mixes different primitive conductors: it has zeros off
$\Re s=3/4$ and infinitely many zeros in $\Re s>1$.  The hypothesis refers
to the final primitive expansion, so terms which cancel must first be
combined.  Merely observing nonmultiplicative coefficients is insufficient,
since a single Dirichlet $L$-function multiplied by a Dirichlet polynomial
can also have nonmultiplicative coefficients.

\subsection{The full lambency-six calculation}
For $m=6$ put
\[
 A=S_{6,1}+S_{6,5},\qquad B=S_{6,1}-S_{6,5},\qquad z=2s-1.
\]
Taking positive and negative integers in their residue classes gives
\begin{align}
 A(\tau)&=\sum_{\substack{n\geq1\\(n,6)=1}}\chi_{-4}(n)nq^{n^2/24},&
 L(A,s)&=24^s(1+3^{-z})L(z,\chi_{-4}),\label{eq:l6-A}\\
 B(\tau)&=\sum_{\substack{n\geq1\\(n,6)=1}}\chi_{-3}(n)nq^{n^2/24},&
 L(B,s)&=24^s(1+2^{-z})L(z,\chi_{-3}).\label{eq:l6-B}
\end{align}
The remaining components are
\begin{align}
 S_{6,2}(\tau)&=2\sum_{\substack{k\geq1\\(k,6)=1}}
 \chi_{-3}(k)kq^{k^2/6},&
 L(S_{6,2},s)&=2\,6^s(1+2^{-z})L(z,\chi_{-3}),\label{eq:l6-2}\\
 S_{6,3}(\tau)&=3\sum_{k\geq1}\chi_{-4}(k)kq^{3k^2/8},&
 L(S_{6,3},s)&=3(8/3)^sL(z,\chi_{-4}),\label{eq:l6-3}\\
 S_{6,4}(\tau)&=4\sum_{k\geq1}\chi_{-3}(k)kq^{2k^2/3},&
 L(S_{6,4},s)&=4(3/2)^sL(z,\chi_{-3}).\label{eq:l6-4}
\end{align}
For example, in $S_{6,2}$ write $n=2k$ and distinguish $k\equiv1,5\pmod6$;
the second residue class comes from a negative $n$.  The other identities
follow in exactly the same way, by the displayed congruence classes.

\begin{theorem}[The first mixed umbral component]\label{thm:l6-mixed-zeros}
For $X=A_5^4D_4$, the first untwined shadow component has transform
\begin{equation}\label{eq:l6-mixed-L}
 L(S^X_1,s)=24^s\left\{
 3(1+3^{1-2s})L(2s-1,\chi_{-4})+
 2(1+2^{1-2s})L(2s-1,\chi_{-3})\right\}.
\end{equation}
It has infinitely many zeros off $\Re s=3/4$, including infinitely many
zeros in $\Re s>1$, with the stripwise linear bounds of
Theorem~\ref{thm:shadow-zero-classification}.
The same conclusions hold for $S^X_5$.
\end{theorem}
\begin{proof}
For this root system, the defining ADE matrices give
$S^X_1=5S_{6,1}+S_{6,5}=3A+2B$ and
$S^X_5=S_{6,1}+5S_{6,5}=3A-2B$.
Substitution of \eqref{eq:l6-A}--\eqref{eq:l6-B} proves the transforms.
Both primitive terms are nonzero, and their conductors are $4$ and $3$,
so Theorem~\ref{thm:shadow-zero-classification} applies.
Alternatively, for the first component its periodic coefficients on
$1,5,7,11\pmod{12}$ are $5,1,-1,-5$.
If $D_c=P L(\cdot,\chi)$ with $P=\sum_{d\leq M}b_dd^{-z}$, then coefficient
uniqueness gives $b_1=5$ and $c(p)=5\chi(p)$ at every prime $p>M$ not dividing
the conductor of $\chi$.  There are arbitrarily large primes
$p\equiv5\pmod{12}$, and there $|c(p)|=1$, a contradiction.
\end{proof}

\subsection{Every lambency-six twining and its symmetric powers}
Let $G=\mathrm{GL}_2(\mathbb F_3)$.  Write $a(g)$ for the number of fixed
lines in $\mathbb F_3^2$, and $b(g)$ for the number of those lines on which
$g$ acts by $+1$ minus the number on which it acts by $-1$.  On the three
partitions of these four lines into two unordered pairs, put
$c(g)=\#\mathrm{Fix}(g)-1$, and let $\delta(g)$ be the sign of this
permutation.  Equivalently $\delta(g)=\det_{\mathbb F_3}(g)\in\{1,-1\}$.
The shadow definition of Cheng, Duncan and Harvey
\cite[Sect.~5.1, equations (5.1)--(5.6), and Table~28, p.~95]{CDH2014} gives
\begin{equation}\label{eq:l6-full-shadow}
 S^X_g=
 \left( (a+1)S_{6,1}+\delta S_{6,5},\;
 bS_{6,2},\;(a+c)S_{6,3},\;bS_{6,4},\;
 \delta S_{6,1}+(a+1)S_{6,5}\right).
\end{equation}
One may check this directly: the four $A_5$ diagrams give $a$ on odd
indices and $b$ on even indices; the $D_4$ triality diagram gives $1$ on
indices $1,5$, $c$ on index $3$, and the two off-diagonal entries
$\delta$.  Its two entries on the even indices cancel because
$S_{6,-r}=-S_{6,r}$.

\begin{theorem}[Symmetric powers in the twined shadow]\label{thm:l6-sympowers}
Choose the faithful two-dimensional complex representation $V$ of $G$
whose trace on one of the order-eight classes is $\sqrt{-2}$.  Then
\[
 \det V=\delta,\qquad
 a-1=\operatorname{tr}\operatorname{Sym}^2V,\qquad
 b=\operatorname{tr}\operatorname{Sym}^3V,\qquad
 a+c=1+\operatorname{tr}\operatorname{Sym}^4V.
\]
Moreover
\[
 V\otimes\operatorname{Sym}^2V
   \simeq\operatorname{Sym}^3V\oplus(V\otimes\det V),\qquad
 V\otimes\det V\simeq\overline V.
\]
The centre $-I$ acts on $\operatorname{Sym}^kV$ by $(-1)^k$.
The nonzero isotypic combinations of every fixed twining are constant
multiples of the five Dirichlet constituents in
\eqref{eq:l6-A}--\eqref{eq:l6-4}.
\end{theorem}
\begin{proof}
Here is the complete class calculation; the two order-eight classes are
combined only in the columns which have equal values.  The faithful
character is the two-dimensional character in the character table of
$\mathrm{GL}_2(3)$ \cite[Table~15, p.~91]{CDH2014}.
\begin{center}
\begin{tabular}{c|rrrrrrr}
class&$1A$&$2A$&$2B$&$4A$&$3A$&$6A$&$8A,8B$\\\hline
size&1&1&12&6&8&8&6,6\\
$\operatorname{tr}V$&2&$-2$&0&0&$-1$&1&$\pm\sqrt{-2}$\\
$\delta$&1&1&$-1$&1&1&1&$-1$\\
$a$&4&4&2&0&1&1&0\\
$b$&4&$-4$&0&0&1&$-1$&0\\
$c$&2&2&0&2&$-1$&$-1$&0
\end{tabular}
\end{center}
The last three rows can also be obtained without a character table by acting
on the four lines and their three pair partitions.  For trace $t$ and
determinant $\delta$ the symmetric-power traces are
\[
 1,\quad t,\quad t^2-\delta,\quad t^3-2\delta t,
 \quad t^4-3\delta t^2+\delta^2.
\]
Substitution in the displayed table proves every character equality.
For completeness, the tensor identity has a canonical proof independent of
this table.  Symmetrization maps $V\otimes\operatorname{Sym}^2V$ onto
$\operatorname{Sym}^3V$; its kernel is the copy of
$V\otimes\bigwedge^2V$ generated by
$v\otimes(uw)-u\otimes(vw)$.  Complete reducibility gives the direct sum.
Multiplication of the trace row by $\delta$ gives its complex conjugate,
proving the second isomorphism.  Faithfulness implies that the nontrivial
central involution acts by $-1$ on $V$, and hence by $(-1)^k$ on its
$k$th symmetric power.
Finally \eqref{eq:l6-full-shadow} gives
\begin{align*}
 S^X_{g,1}+S^X_{g,5}&=(a+1+\delta)A,\\
 S^X_{g,1}-S^X_{g,5}&=(a+1-\delta)B,
\end{align*}
while the other three entries are $bS_{6,2}$, $(a+c)S_{6,3}$, and
$bS_{6,4}$.  These are exactly the asserted constituents.
\end{proof}

The zero cases must be excluded when making any Riemann-hypothesis
statement.  Explicitly the respective multipliers of
$(A,B,S_{6,2},S_{6,3},S_{6,4})$, where the first two mean the sum and
difference of the twined components, are
\begin{center}
\begin{tabular}{c|rrrrr}
class&$A$&$B$&$S_{6,2}$&$S_{6,3}$&$S_{6,4}$\\\hline
$1A$&6&4&4&6&4\\
$2A$&6&4&$-4$&6&$-4$\\
$2B$&2&4&0&2&0\\
$4A$&2&0&0&2&0\\
$3A$&3&1&1&0&1\\
$6A$&3&1&$-1$&0&$-1$\\
$8A,8B$&0&2&0&0&0
\end{tabular}
\end{center}
Every nonzero entry inherits the zero line of its constituent.  A zero entry
is the identically zero function and has no discrete zero-counting problem.
For a single first component, as opposed to its sum and difference with the
fifth, the coefficients of $A$ and $B$ are
$(a+1+\delta)/2$ and $(a+1-\delta)/2$.
Thus the first and fifth components have the unconditional mixed-character
zero phenomenon on $1A,2A,2B,3A,6A$; they are single-character constituents
on $4A,8A,8B$.

If an elliptic curve $E/\mathbb Q$ has full mod-$3$ image, the Weil pairing
identifies $\det\overline\rho_{E,3}$ with the mod-$3$ cyclotomic character.
Consequently, at an unramified Frobenius,
$\delta(\overline\rho_{E,3}(\mathrm{Frob}_p))=\chi_{-3}(p)$.
The remaining character rows are traces of the finite-image Artin
representations $\operatorname{Sym}^j(V\circ\overline\rho_{E,3})$.
Their Artin $L$-functions are different objects from the unary transforms
of a fixed twining; the Dirichlet theorem directly controls the determinant
character but supplies no Riemann hypothesis for these higher-dimensional
Artin representations.

\subsection{The elementary lattice factorizations}
The following formulas use $(v,v)/2$ as the quadratic form for $A_2,D_4,E_8$,
and the ordinary Gaussian norm for $\mathbb Z[i]$:
\begin{align}
 Z_{A_2}(s)&=6\zeta(s)L(s,\chi_{-3}),&
 Z_{\mathbb Z[i]}(s)&=4\zeta(s)L(s,\chi_{-4}),\label{eq:CM-Epstein}\\
 Z_{D_4}(s)&=24(1-2^{1-s})\zeta(s)\zeta(s-1),&
 Z_{E_8}(s)&=240\zeta(s)\zeta(s-3).\label{eq:DE-Epstein}
\end{align}
For the first pair, $\mathbb Z[\omega]$ and $\mathbb Z[i]$ are principal
ideal domains with respectively six and four units.  Grouping nonzero
elements by their principal ideals gives the respective Dedekind zeta
functions, which factor as displayed.  For $D_4$, use
$D_4=\{x\in\mathbb Z^4:\sum x_i\equiv0\pmod2\}$.
The condition $\sum x_i^2=2n$ already imposes this parity condition, so
Jacobi's four-square formula gives
\[
 r_{D_4}(n)=r_4(2n)=8\sum_{d\mid2n,\,4\nmid d}d
       =24\sum_{d\mid n,\,d\ {
m odd}}d.
\]
Taking its Dirichlet series proves the third identity.  The theta series of
$E_8$ is the unique normalized weight-four modular form on
$\mathrm{SL}_2(\mathbb Z)$, namely $E_4=1+240\sum_{n\geq1}\sigma_3(n)q^n$;
this proves the fourth.

Thus the Dirichlet Riemann hypothesis places the nontrivial Dirichlet-factor
zeros on $\Re s=1/2$ for the first pair, on $\Re s=1/2,3/2$ for $D_4$,
and on $\Re s=1/2,7/2$ for $E_8$.  The factor $1-2^{1-s}$ contributes
zeros on $\Re s=1$, except that its zero at $s=1$ cancels the pole of
$\zeta(s)$ and is not a zero of $Z_{D_4}$.
Similarly, the signed transform
\[
 (4^s-2)\,6\zeta(s)L(s,\chi_{-3})
\]
has all its zeros in $0<\Re s<1$ on $\Re s=1/2$, since the additional
factor vanishes precisely at
$s=\tfrac12+\pi i k/\log2$, $k\in\mathbb Z$.
At lambency two, Jacobi's identity
$\eta(\tau)^3=\sum_{n\geq1}\chi_{-4}(n)nq^{n^2/8}$ gives
$L(24\eta^3,s)=24\,8^sL(2s-1,\chi_{-4})$.

\subsection{The scalar components at all twenty-three lambencies}
\label{sec:all-lambency-classification}
The character criterion gives a finite, exact classification of every
untwined component.  We take the independent indices $1\leq r<m$; the
remaining indices are their negatives or are identically zero.

\begin{proposition}[All scalar umbral shadows]\label{prop:all-scalar-shadows}
Among the $303$ components $S^X_r$, with $X$ running through the twenty-three
Niemeier root systems and $1\leq r<m(X)$, exactly $92$ vanish identically,
$39$ have one primitive Dirichlet character in
\eqref{eq:primitive-shadow-expansion}, and $172$ have at least two.
Every one of those $172$ components has the unconditional off-line zeros
of Theorem~\ref{thm:shadow-zero-classification}.
The single-character and zero components are specified in the table below;
all components not listed in either category are mixed.

Of the $39$ single-character components, $21$ have a monomial
$P_\psi(z)$, ten have a monomial times $1+p^{-z}$, and eight have a monomial
times $1+p^{1-z}$, with $p=2$ or $3$.
For these last eight, the additional zeros are
\begin{equation}\label{eq:scalar-shadow-boundary-zeros}
 s=1+\frac{(2j+1)\pi i}{2\log p},\qquad j\in\mathbb Z.
\end{equation}
They are present unconditionally and cannot be cancelled by a Dirichlet
$L$-function, which is nonvanishing at $\Re(2s-1)=1$.
The eight cases are
\[
 \begin{array}{c|c|c}
 X&r&p\\\hline
 A_{11}D_7E_6&4,8&2\\
 E_6^4&4,8&2\\
 A_{17}E_7&9&3\\
 D_{10}E_7^2&3,9,15&3.
 \end{array}
\]
Consequently, under the Dirichlet Riemann hypothesis, every zero of any of
the $39$ single-character components with $\Re s>1/2$ lies on
$\Re s=3/4$ or on the explicitly described line $\Re s=1$.
\end{proposition}

\begin{center}
\small
\begin{tabular}{l r l r r}
$X$&$m$&single-character indices&zero count&mixed count\\\hline
$A_1^{24}$&2&1&0&0\\
$A_2^{12}$&3&1,2&0&0\\
$A_3^8$&4&2&0&2\\
$A_4^6$&5&none&0&4\\
$A_5^4D_4$&6&2,3,4&0&2\\
$A_6^4$&7&none&0&6\\
$A_7^2D_5^2$&8&2,4,6&0&4\\
$A_8^3$&9&3,6&0&6\\
$A_9^2D_6$&10&5&0&8\\
$A_{11}D_7E_6$&12&2,4,6,8,10&0&6\\
$A_{12}^2$&13&none&0&12\\
$A_{15}D_9$&16&4,8,12&0&12\\
$A_{17}E_7$&18&6,9,12&0&14\\
$A_{24}$&25&none&0&24\\
$D_4^6$&6&1,3,5&2&0\\
$D_6^4$&10&5&4&4\\
$D_8^3$&14&7&6&6\\
$D_{10}E_7^2$&18&3,9,15&8&6\\
$D_{12}^2$&22&11&10&10\\
$D_{16}E_8$&30&5,15,25&14&12\\
$D_{24}$&46&23&22&22\\
$E_6^4$&12&4,8&5&4\\
$E_8^3$&30&none&21&8\\\hline
Total&&39&92&172
\end{tabular}
\end{center}
For the seven root systems beginning with $D$, precisely the even indices
vanish.  For $E_6^4$ the zero indices are $2,3,6,9,10$.
For $E_8^3$ the nonzero indices are $1,7,11,13,17,19,23,29$.
Thus the table specifies every index, not only the counts.

\begin{proof}
We describe a finite certificate using exact arithmetic.  For $d\mid m$,
the defining matrix is
\[
 \Omega_m(d)_{r,t}=\begin{cases}
 1,&r+t\equiv0\pmod{2d},\quad r-t\equiv0\pmod{2m/d},\\
 0,&\text{otherwise}.
 \end{cases}
\]
The ADE matrices are, respectively,
\[
 \begin{array}{c|c}
 A_{m-1}&\Omega_m(1)\\
 D_{m/2+1}&\Omega_m(1)+\Omega_m(m/2)\\
 E_6&\Omega_{12}(1)+\Omega_{12}(4)+\Omega_{12}(6)\\
 E_7&\Omega_{18}(1)+\Omega_{18}(6)+\Omega_{18}(9)\\
 E_8&\Omega_{30}(1)+\Omega_{30}(6)+\Omega_{30}(10)+\Omega_{30}(15).
 \end{array}
\]
These are Cheng--Duncan--Harvey's defining matrices
\cite[Table~5 and Sect.~3.3]{CDH2014}; a union of root systems means
adding their matrices.  For each row put
$c_r(n)=\Omega^X_{r,n}-\Omega^X_{r,-n}$, and apply
\eqref{eq:character-stratum-coefficients} followed by
\eqref{eq:imprimitive-shadow}.

The computation enumerates the characters using the cyclic unit groups of
odd prime powers and
$(\mathbb Z/2^a\mathbb Z)^\times=\langle-1\rangle\times\langle5\rangle$
for $a\geq3$.  All sums are evaluated in the exact ring
$\mathbb Q[x]/(\Phi_e(x))$, where $e$ is the exponent of the relevant unit
group.  The conductor is the smallest divisor $f$ of the modulus for which
the character is trivial on the kernel of reduction modulo $f$.
After collecting equal primitive characters, every resulting polynomial
is expanded back into its coefficient sequence.  Equality with $c_r$ is
checked at every residue of its full period $2m$; the reconstructed terms
have periods dividing $2m$, so this check proves each identity.
The complete tables and the program \emph{Shadow character classification}
are in the workbench.

The table then follows by counting the nonzero primitive constituents.
For the single-character cases the exact polynomials are scalar multiples
of $d^{-z}$, $d^{-z}(1+p^{-z})$, or $d^{-z}(1+p^{1-z})$.
To show the last possibility explicitly, the two components with
$X=A_{11}D_7E_6$, $r=4,8$, have
\[
 D_{c_r}(z)=2\,4^{-z}(1+2^{1-z})L(z,\chi_{-3}),
\]
and the two with $X=E_6^4$, $r=4,8$, have twice this expression.
The component $X=A_{17}E_7$, $r=9$, has
\[
 D_{c_9}(z)=3^{-z}(1+3^{1-z})L(z,\chi_{-4});
\]
the components $r=3,15$ of $D_{10}E_7^2$ have this same expression,
and its component $r=9$ has twice it.  Solving $p^{1-z}=-1$ with
$z=2s-1$ gives \eqref{eq:scalar-shadow-boundary-zeros}.
The zero theorem for the other $172$ components is
Theorem~\ref{thm:shadow-zero-classification}.
\end{proof}

\subsection{The eta series and the signed hexagonal Gaussian}
\label{sec:eta-signed-mellin}
Let $\chi_{12}(n)=(12/n)$ be the primitive even character of conductor
$12$, with values $1,-1,-1,1$ at $1,5,7,11$ modulo $12$.
Euler's pentagonal identity gives the exact unary expansion
\begin{equation}\label{eq:eta-unary-chi12}
 \eta(\tau)=\sum_{n\geq1}\chi_{12}(n)q^{n^2/24}.
\end{equation}
Indeed its pentagonal exponent is
$1/24+k(3k-1)/2=(6k-1)^2/24$, and taking the absolute values of
$6k-1$, for $k\in\mathbb Z$, enumerates the positive integers prime to
$6$ exactly once with sign $\chi_{12}$.
Thus, with the Mellin Dirichlet convention that a term $a_\nu q^\nu$
contributes $a_\nu\nu^{-s}$,
\begin{equation}\label{eq:eta-mellin-dirichlet}
 L(\eta,s)=24^sL(2s,\chi_{12}),\qquad
 \int_0^\infty\eta(iy)y^{s-1}\,dy
   =\Gamma(s)\left(\frac{12}{\pi}\right)^s L(2s,\chi_{12}).
\end{equation}
Termwise integration is initially valid for $\Re s>1/2$.
The transformation of $\eta$ shows exponential decay, up to a power,
at both integration endpoints, so the completed integral extends
entirely.  If the sum in \eqref{eq:eta-unary-chi12} were instead taken over
all $n\in\mathbb Z$, it would need a factor $1/2$.
Alp\"oge's theorem places the nontrivial zeros of $L(\eta,s)$
on $\Re s=1/4$.

For the signed Gaussian, use the Euclidean hexagonal lattice
$L=\mathbb Z+\omega\mathbb Z$ with
$\omega=(-1+\sqrt{-3})/2$ and $|a+b\omega|^2=a^2-ab+b^2$.
There are four cosets of $L$ in $\frac12L$.  Define
\[
 \epsilon(v)=\begin{cases}-1,&v\in L,\\+1,&v\in\frac12L\setminus L,
 \end{cases}
 \qquad
 G_\epsilon(t)=\sum_{v\in\frac12L}\epsilon(v)e^{-\pi t|v|^2}.
\]
The convention fixes the lattice scale: the quadratic form is $|v|^2$,
without division by $\Im\omega$.  Writing
$\Theta_L(t)=\sum_{v\in L}e^{-\pi t|v|^2}$ gives
\[
 G_\epsilon(t)=\Theta_L(t/4)-2\Theta_L(t),\qquad G_\epsilon(\infty)=-1.
\]
For $\Re s>1$, absolute convergence and termwise integration yield
\begin{align}
 \frac{\pi^s}{\Gamma(s)}
   \int_0^\infty\bigl(G_\epsilon(t)+1\bigr)t^{s-1}\,dt
 &=\sum_{0\ne v\in\frac12L}\epsilon(v)|v|^{-2s}\notag\\
 &=(4^s-2)Z_L(s)\notag\\
 &=6(4^s-2)\zeta(s)L(s,\chi_{-3}).
 \label{eq:signed-hexagonal-mellin}
\end{align}
The last equality follows by identifying nonzero Eisenstein integers
modulo their six units with nonzero ideals; the Dedekind zeta function of
$\mathbb Q(\sqrt{-3})$ is $\zeta(s)L(s,\chi_{-3})$.
Both sides then continue meromorphically.
The extra elementary factor has zeros
\[
 s=\frac12+\frac{\pi i k}{\log2},\qquad k\in\mathbb Z.
\]
If the area-normalized quadratic form $|v|^2/\Im\omega$ is used instead,
the right side acquires the nonvanishing factor $(\Im\omega)^s$.

\section{Off-line zeros of every Niemeier Epstein zeta function}
\label{sec:niemeier-zeros}
The cusp contribution to a rank-twenty-four theta series does more than
prevent a direct application of the Dirichlet Riemann hypothesis: it
forces infinitely many zeros far from the Epstein symmetry line.
The required analytic result is the hybrid-universality theorem of
Nakamura and Pa\'nkowski \cite[Main Theorems~1--2]{NP2016}.

Let $L$ be any even unimodular positive-definite lattice of rank $24$,
including the Leech lattice, and put
\[
 \Theta_L(\tau)=\sum_{v\in L}q^{(v,v)/2},\qquad
 N_2(L)=\#\{v\in L:(v,v)=2\},\qquad A=\frac{65520}{691}.
\]
Since $M_{12}(\mathrm{SL}_2(\mathbb Z))$ is spanned by $E_{12}$ and
$\Delta$, comparison of the constant and first Fourier coefficients gives
\begin{equation}\label{eq:niemeier-zeta-decomposition}
 \Theta_L=E_{12}+\bigl(N_2(L)-A\bigr)\Delta,\qquad
 Z_L(w)=A\zeta(w)\zeta(w-11)+\bigl(N_2(L)-A\bigr)L(\Delta,w).
\end{equation}
Here $Z_L(w)=\sum_{v\ne0}((v,v)/2)^{-w}$ initially for $\Re w>12$.

\begin{theorem}[Niemeier zeros]
\label{thm:niemeier-offline}
For every even unimodular positive-definite rank-$24$ lattice $L$ and every
$23/2<\alpha<\beta<12$, there are positive constants $c_1,c_2,T_0$ such that
for $T\geq T_0$,
\[
 c_1T\leq\#\{w:Z_L(w)=0,\ \alpha<\Re w<\beta,\ 0<\Im w<T\}
       \leq c_2T.
\]
The upper bound counts multiplicities; the lower bound can count distinct
simple zeros.  The same assertion holds in the reflected strip
$12-\beta<\Re w<12-\alpha$.
In particular the Leech Epstein zeta function has infinitely many zeros
off $\Re w=6$, and so does every Niemeier Epstein zeta function.
No Riemann hypothesis is an assumption of this theorem.
\end{theorem}
\begin{proof}
Set $z=w-11$, $B=N_2(L)-A$, and consider the polynomial in $X$
\begin{equation}\label{eq:niemeier-hybrid-polynomial}
 P_z(X)=A\zeta(z+11)X+B L(\Delta,z+11).
\end{equation}
Both coefficients are Dirichlet series absolutely convergent for
$\Re z>1/2$.  For the second coefficient even the elementary bound
$\tau(n)=O(n^6)$ suffices: boundedness of
$y^6|\Delta(x+iy)|$ on the modular fundamental domain, and hence on the
upper half-plane, combined with Fourier inversion at $y=1/n$ proves that
bound.  Therefore $L(\Delta,z+11)$ converges absolutely already for
$\Re z>-4$.  The first coefficient is even more plainly absolutely
convergent in the required half-plane.

Neither coefficient is identically zero.  Indeed $A\ne0$, the two
$L$-series tend to $1$ as the real part tends to infinity, and
$B\ne0$ because $N_2(L)$ is an integer whereas $65520/691$ is not.
Thus $P_z$ is a nonmonomial polynomial of positive degree.
The Riemann zeta function is hybridly universal in
$1/2<\Re z<1$.  Main Theorem~1 of Nakamura--Pa\'nkowski
\cite[arXiv v4, p.~4]{NP2016} now supplies a lower bound proportional to
$T$ for the zeros of
\[
 P_z(\zeta(z))=Z_L(z+11)
\]
in every closed interior vertical strip of that region.
Their Main Theorem~2 gives the corresponding upper bound: its hypotheses
for $\zeta(z)$ are its standard meromorphic continuation, finite order,
and $\int_0^T|\zeta(\sigma+it)|^2dt=O_\sigma(T)$ for fixed
$1/2<\sigma<1$.

For clarity, the lower bound may be made a bound for simple zeros directly
from the same hybrid theorem.  Choose a small closed disc $K$ about a point
$z_0$ in the desired strip where the constant coefficient
$B L(\Delta,z+11)$ is nonzero.  Its existence follows because that
coefficient is a nonzero analytic function.  Shrink the disc so that it is
nonzero throughout $K$.  The coefficient $A\zeta(z+11)$ is nonzero there
by the Euler product.  For a sufficiently small nonzero $\epsilon$, the
function
\[
 f(z)=\frac{\epsilon(z-z_0)-B L(\Delta,z+11)}{A\zeta(z+11)}
\]
is analytic and nonvanishing on $K$, and
$P_z(f(z))=\epsilon(z-z_0)$ has exactly one simple zero inside $K$.
Theorem~E of Nakamura--Pa\'nkowski \cite[arXiv v4, p.~5]{NP2016} gives a
set of shifts $t\in[0,T]$ of measure at least $cT$ for which
$P_{z+it}(\zeta(z+it))$ uniformly approximates $P_z(f(z))$ on $K$.
Choose the error smaller than $|\epsilon|$ times the radius of $K$.
Rouch\'e's theorem then gives exactly one zero, counted with multiplicity,
in every translated disc, hence a simple zero.  One zero can account for
at most an interval of shifts of length twice that radius.  There are
therefore at least a positive multiple of $T$ distinct simple zeros.
Translating $z$ to $w=z+11$ proves the first assertion.

Finally Poisson summation gives
\[
 (2\pi)^{-w}\Gamma(w)Z_L(w)
  =(2\pi)^{-(12-w)}\Gamma(12-w)Z_L(12-w).
\]
Both gamma factors are finite and nonzero in the two strips in question.
Since the coefficients of $Z_L$ are real, reflection sends a zero $w$ in
the upper half-plane to $12-\overline w$, also in the upper half-plane,
with the same multiplicity.  This proves the reflected assertion.
\end{proof}

For the Leech lattice $N_2=0$, so the polynomial in
\eqref{eq:niemeier-hybrid-polynomial} is
$A\{\zeta(z+11)X-L(\Delta,z+11)\}$.
The theorem therefore applies to the Leech lattice without a separate
numerical search.  It locates zeros in every substrip of
$11.5<\Re w<12$ and $0<\Re w<0.5$; it does not claim that these regions
contain all zeros, or decide whether there are zeros in $\Re w>12$.
The rank-two class-number theorem of Davenport and Heilbronn cannot simply
be applied to a rank-$24$ lattice.  The proof above instead verifies the
hypotheses of a theorem that applies to the precise rank-$24$ decomposition.

\section{The residue action and its four eta products}
This section identifies the modular forms arising from the cyclic action on
four copies of the root lattice $A_5$.
Let $N$ be the Niemeier lattice with root system $A_5^4D_4$.
Let $c$ cyclically permute the six coordinates of each $A_5$ factor and
act trivially on $D_4$.  Root reflections act trivially on the
root-lattice discriminant group, so $c$ preserves the glue defining $N$.
A Frame shape $\prod_d d^{a_d}$ means the identity
$\det(t-c)=\prod_d(t^d-1)^{a_d}$.

\begin{theorem}\label{thm:residue-eta}
The squares in $(\mathbb Z/840\mathbb Z)^\times$ form the cyclic group
$S=\langle289\rangle$ of order six. Under $289\mapsto c$, an element
of order $d$ has Frame shape $d^{24/d}$. Its eta product is
$F_d(\tau)=\eta(d\tau)^{24/d}$, and the four products are normalized
newforms, as follows:
\[
\begin{array}{c|c|c|c}
 d&F_d&\text{weight}&\text{level}\\\hline
 1&\Delta&12&1\\
 2&\eta(2\tau)^{12}&6&4\\
 3&\eta(3\tau)^8&4&9\\
 6&\eta(6\tau)^4&2&36
\end{array}
\]
Exactly the last two have complex multiplication, in both cases by
$\mathbb Q(\sqrt{-3})$. This occurs exactly when the cubic character
modulo seven is nontrivial on the element of $S$.
\end{theorem}
\begin{proof}
The Chinese remainder theorem gives the group of squares modulo
$840=8\cdot3\cdot5\cdot7$ as the product of groups of orders
$1,1,2,3$. The successive powers of $289$ are
$1,289,361,169,121,529$, and the sixth power is one. The action of a
six-cycle on the sum-zero subspace of $\mathbb Q^6$ has characteristic
polynomial $(t^6-1)/(t-1)$. Four copies, together with the four fixed
coordinates from $D_4$, therefore have polynomial $(t^6-1)^4$.
Equivalently, the rational representation is four copies of the regular
representation of the cyclic group of order six. Restriction to an
element of order $d$ gives $24/d$ regular representations of order $d$.

For $N=d^2$ the eta-product congruences are
$\sum_{\delta\mid N}\delta r_\delta=24$ and
$\sum_{\delta\mid N}(N/\delta)r_\delta=24$, with $r_d=24/d$.
The eta transformation formula gives a modular form with trivial
character; all cusp orders are positive because its only eta exponent
is positive. Its first term is $q$. The spaces
$S_{12}(\mathrm{SL}_2(\mathbb Z))$, $S_6(\Gamma_0(4))$,
$S_4(\Gamma_0(9))$, and $S_2(\Gamma_0(36))$ each have dimension one.
Their lower-level old spaces are zero, so these normalized cusp forms
are newforms and simultaneous Hecke eigenforms.

The complex multiplication assertions for $d=3,6$ also follow from the
explicit characters in Theorem~\ref{thm:primitive-cube}.
A complex multiplication field for a newform has discriminant dividing
its level. No imaginary quadratic discriminant divides one. At level
four the only possibility is $\mathbb Q(i)$, whereas the coefficient
of $q^3$ in $\eta(2\tau)^{12}$ is $-12$; it would vanish at the inert
prime three for a form with this complex multiplication. Thus the first
two forms have no complex multiplication. Finally, if $r=289^j$, its
cubic character modulo seven is trivial exactly when $3\mid j$.
Those are precisely the elements of orders one and two.
\end{proof}

\begin{proposition}[The two residue symbols]\label{prop:residue-symbol}
Put $\omega=e^{2\pi i/3}$, choose the quartic character modulo five
by $(2/5)_4=i$, and put
$\pi=3+\omega^2=2-\omega$, a prime of norm seven in $\mathbb Z[\omega]$.
Then the character of $S$ which sends $289$ to $e^{\pi i/3}$ is
\[
 \kappa(r)=\left(\frac r5\right)_4
                \left(\frac r\pi\right)_3 .
\]
In particular its cubic factor is the restriction to rational integers
of the cubic residue character at $\pi$.
\end{proposition}
\begin{proof}
In $\mathbb Z[\omega]/(\pi)=\mathbb F_7$ one has $\omega=2$,
so $(r/\pi)_3$ is the unique cube root of unity congruent to
$r^{(7-1)/3}=r^2$ modulo $\pi$. At $r=289$ the quartic factor is
$-1$, while $r=2$ modulo seven gives the cubic factor $\omega^2$.
Their product is $-\omega^2=e^{\pi i/3}$. Both factors are
characters; equality at the generator proves the identity on $S$.
\end{proof}

\section{The primitive cubic character and all local factors}
We distinguish a power of a primitive Hecke character from the
imprimitive ideal series obtained by retaining its old excluded primes.
Put $\omega=e^{2\pi i/3}$, $\mathcal O=\mathbb Z[\omega]$, and
$\mathfrak p=(\sqrt{-3})$. The ring $\mathcal O$ is a principal ideal
domain. Every ideal prime to $2\mathfrak p$ has a unique generator
$\alpha\equiv1\pmod{2\mathfrak p}$; every ideal prime to $3$ has a
unique generator $\beta\equiv1\pmod3$. Define their respective
Hecke characters by $\psi_{36}(\mathfrak a)=\alpha$ and
$\psi_{27}(\mathfrak a)=\beta$.

\begin{theorem}\label{thm:primitive-cube}
The characters $\psi_{36}$ and $\psi_{27}$ have infinity type one
and conductors $2\mathfrak p$ and $(3)$. Their primitive cubes coincide:
\[
 \psi_{36}^3=\psi_{27}^3=\Psi,
 \qquad \mathfrak f_\Psi=\mathfrak p.
\]
For an ideal prime to $3$, choose any generator
$\gamma\equiv1\pmod{\mathfrak p}$. Then
$\Psi(\mathfrak a)=\gamma^3$ is well-defined, and
\[
 f=\theta_{\psi_{36}}=\eta(6\tau)^4,
 \quad g=\theta_\Psi=\eta(3\tau)^8,
 \quad \theta_{\psi_{27}}=\eta(3\tau)^2\eta(9\tau)^2.
\]
For the usual arithmetic normalization of the finite local factors,
\begin{equation}\label{eq:full-cube}
 L(\operatorname{Sym}^3f,s)=L(g,s)L(f,s-1)
\end{equation}
with no omitted primes. At $p=2$ the factor on either side is
$(1+8\,2^{-2s})^{-1}$; at $p=3$ it is one. At every prime $p\ge5$,
\[
 a_p(g)=a_p(f)^3-3p\,a_p(f).
\]
\end{theorem}
\begin{proof}
Reduction identifies $\mathcal O^\times$ with
$(\mathcal O/(2\mathfrak p))^\times$ and with
$(\mathcal O/(3))^\times$: all three groups have order six, and the
six units have distinct reductions. This proves the two uniqueness
assertions and multiplicativity. Modulo $\mathfrak p$, the units with
value one are exactly $1,\omega,\omega^2$. Their cubes are one,
so the displayed definition of $\Psi$ is independent of the
choice. The primary generators for both type-one characters are
primary modulo $\mathfrak p$, proving the cube identities away from
six. They therefore identify the primitive characters.

The finite part of a type-one character cannot be trivial on any unit
other than one. Removing either prime factor from $2\mathfrak p$
would make $-1$ or $\omega$ congruent to one; replacing $(3)$ by
$\mathfrak p$ would make $\omega$ congruent to one. Their conductors
are therefore exactly as stated. For type three, modulus one is
impossible because $(-1)^3=-1$, while modulus $\mathfrak p$ works.
The theta identities are the explicit formulas of
Huber, Liu, McLaughlin, Ye, Yuan and Zhang
\cite[Lemma~2.1, equations~(2.2), (2.4), and~(2.10)]{HuberCM}.
In the last formula
each ideal has three primary generators modulo $\mathfrak p$, which
accounts for its factor $1/3$.

Let $\rho=\operatorname{Ind}_{G_K}^{G_{\mathbb Q}}\psi_{36,\ell}$
be the compatible two-dimensional representation of $f$, and let
$\varepsilon$ be the cyclotomic character. Its determinant is
$\varepsilon$. On restriction to $G_K$, write a basis as $e_1,e_2$
with characters $\psi,\psi^c$. In the symmetric cube, the spans
$\langle e_1^3,e_2^3\rangle$ and
$\langle e_1^2e_2,e_1e_2^2\rangle$ are invariant. The first is
$\operatorname{Ind}\psi^3$. The second restricts to
$\psi\varepsilon\oplus\psi^c\varepsilon$ and is isomorphic to
$\rho\otimes\varepsilon$; an apparent quadratic sign in a choice of
induction basis is removed by the standard isomorphism
$\operatorname{Ind}\psi\simeq
 (\operatorname{Ind}\psi)\otimes\chi_K$.
Thus
\[
 \operatorname{Sym}^3\rho\simeq
 \operatorname{Ind}\Psi\oplus(\rho\otimes\varepsilon).
\]
The local Euler factor is the inverse characteristic polynomial of
Frobenius on inertia invariants. Taking inertia invariants respects
direct sums, so this representation identity proves
\eqref{eq:full-cube} at every finite prime, including ramified primes.

At two, $f$ has trivial local factor, while two is inert in $K$ and
$\Psi((2))=(-2)^3=-8$; its norm is four. This gives the claimed
factor. At three, both characters have ramified local components
and the induced representations have no inertia invariants, so both
factors are one. At a split unramified prime, the two Frobenius
eigenvalues $\alpha,\beta$ satisfy $\alpha\beta=p$, and
$\alpha^3+\beta^3=(\alpha+\beta)^3-3p(\alpha+\beta)$.
At an inert prime both traces vanish. This proves the coefficient
identity in all cases.
\end{proof}

\section{Theta coordinates of ternary perpendicular sublattices}
The same two small shell counts determine the theta series for every
isotropic plane modulo three.
Write $E=E_{12}=1+C\sum_{n\ge1}\sigma_{11}(n)q^n$, where
$C=65520/691$, and put
\[
 a=\frac{7381}{66430},\qquad b=\frac{59049}{66430}.
\]
Let $F\in S_{12}(\Gamma_0(3))$ be the normalized newform whose
first coefficients are $q+78q^2-243q^3+4036q^4+\cdots$.

\begin{lemma}\label{lem:level-three-basis}
Put $H=\eta(\tau)^{12}\eta(3\tau)^{12}$. Then
\[
 F=\Delta(\tau)+729\Delta(3\tau)+102H,
 \qquad F|_{12}W_3=F.
\]
The forms $E(\tau),E(3\tau),\Delta(\tau),\Delta(3\tau),F(\tau)$
are a basis of $M_{12}(\Gamma_0(3))$.
\end{lemma}
\begin{proof}
The eta transformation law shows that $H$ is a cusp form of weight
twelve and trivial character at level three. Direct multiplication gives
$H=q^2-12q^3+54q^4+\cdots$. The index of $\Gamma_0(3)$ is four,
so the valence bound gives an injective map from $M_{12}(\Gamma_0(3))$
to its coefficients of $q^0,\ldots,q^4$. The five displayed forms with
$H$ in place of $F$ have independent images. Thus they are a basis;
the cusp space has dimension three and its old subspace is spanned by
$\Delta(\tau),\Delta(3\tau)$.

With the normalization
$T_2\sum c_nq^n=\sum(c_{2n}+2^{11}c_{n/2})q^n$,
coefficient comparison through $q^4$ gives
\[
 T_2H=\Delta+729\Delta(3\tau)+78H.
\]
Here $c_{n/2}=0$ if $n$ is odd; the finite comparison uses the product
for $H$ through $q^8$. Both oldforms have $T_2$ eigenvalue $-24$.
The displayed linear combination $F$ has eigenvalue $78$, so it
spans the one-dimensional new subspace and is its normalized newform.
Finally $\Delta|W_3=729\Delta(3\tau)$ and $H|W_3=H$, proving
the Fricke sign. These comparisons are reproduced with exact rational
coefficients in the accompanying coefficient program.
\end{proof}

\begin{theorem}\label{thm:universal-theta-coordinates}
Let $\Lambda$ be the Leech lattice and let $P$ be a totally isotropic
two-dimensional subspace of $\Lambda/3\Lambda$, with quadratic form
$x\mapsto(x,x)/2\pmod3$. Set
\[
 J_P=\{x\in\Lambda:(x,P)\equiv0\pmod3\},
\]\[ 
 k=\#\{x\in J_P:(x,x)=4\},\qquad
 t=\#\{u\in3J_P^\vee:(u,u)=6\}.
\]
Then $J_P$ has determinant $81$, level three, and theta series
\begin{equation}\label{eq:universal-theta}
 \Theta_{J_P}=aE(\tau)+bE(3\tau)
            +x\Delta(\tau)+y\Delta(3\tau)+zF(\tau),
\end{equation}
where
\[
 z=\frac{k-2073aC}{102},\qquad
 x=-aC-z,\qquad
 y=59049t-bC-729z.
\]
Moreover $t\in\{0,2,4,6,8\}$: half of $t$ is the number of
lines of $P$ represented by norm-six vectors of $\Lambda$.
\end{theorem}
\begin{proof}
Unimodularity makes the pairing modulo three nondegenerate, so the
two independent conditions defining $J_P$ give index nine.
Consequently its determinant is $9^2$. If $\widetilde P$ denotes
any lifts of a basis of $P$, then
$J_P^\vee=\Lambda+\tfrac13\mathbb Z\widetilde P$.
Total isotropy implies $3J_P^\vee\subset J_P$ and
$3(u,u)/2\in\mathbb Z$ for every $u\in J_P^\vee$.
The associated discriminant form is nondegenerate and nontrivial,
so its level is exactly three. Poisson summation and the evenness of
$J_P$ give modularity of weight twelve and trivial character on
$\Gamma_0(3)$, with Fricke transformation
\[
 \Theta_{J_P}|_{12}W_3
       =81\,\Theta_{J_P^\vee}(3\tau).
\]
Here the determinant-normalized slash operator is used, so that
\[E|W_3=729E(3\tau),\qquad E(3\tau)|W_3=E/729.\]

The space $M_{12}(\Gamma_0(3))$ has basis
\[E(\tau),E(3\tau),\Delta(\tau),\Delta(3\tau),F(\tau).\]
The newform satisfies $F|W_3=F$. Comparing constant terms at the two
cusps gives $a+b=1$ and $729a+b/729=81$, hence the displayed
constants. Since $J_P$ contains no roots, its $q$ coefficient is
zero, giving $x+z=-aC$. Its $q^2$ coefficient gives
$k=2049aC-24x+78z=2073aC+102z$.
The coefficient of $q$ after applying $W_3$ is on the one hand
$81t$, and on the other hand $(bC+y)/729+z$. This yields $y$.

Finally $3J_P^\vee$ is the preimage of $P$ in $\Lambda$.
There is at most one norm-six vector in each nonzero residue class
modulo $3\Lambda$: if $u,v$ were distinct, their squared distance
would be at least $36$, the minimum of $3\Lambda$, but at most
$(\sqrt6+\sqrt6)^2=24$. The eight nonzero vectors of $P$ form
four pairs of opposite vectors. Opposite classes have the same
norm-six count, proving the last assertion.
\end{proof}

\section{Exact theta coordinates for the 129 cases}
The theta-coordinate formula reduces each lattice calculation to two
small shell enumerations. We construct the lattices from their defining
glue and plane representatives, then enumerate the two shells exactly.

\begin{theorem}[Finite computation on the representatives]
\label{thm:129-theta}
For each of the 129 plane representatives, the reconstructed lattice
is even, positive definite, rootless, of determinant $81$ and level
three. The pairs $(k,t)$ in
Theorem~\ref{thm:universal-theta-coordinates} give precisely the
18 theta series in Table~\ref{tab:18theta}. The last column counts
representatives, not orbit sizes. Complete rational values of all five
coordinates, separately for every representative, are tabulated by
the workbench program \emph{Ternary theta coordinates}.
\end{theorem}

\begin{table}[htbp]
\centering\small
\caption{The 18 theta series among the 129 plane representatives.
The two Eisenstein coordinates are always $a,b$; $x,y$ are recovered
from $k,t,z$ by Theorem~\ref{thm:universal-theta-coordinates}.}
\label{tab:18theta}
\begin{tabular}{rrrr}
\hline
$k$&$t$&$z$&Number of representatives\\\hline
21600&0&$-2916/1241$&28\\
21762&0&$-945/1241$&15\\
21924&0&$1026/1241$&23\\
22086&0&$2997/1241$&25\\
22248&0&$4968/1241$&11\\
22410&0&$6939/1241$&4\\
22572&0&$8910/1241$&7\\
22896&0&$756/73$&2\\
23058&0&$14823/1241$&1\\
23544&0&$20736/1241$&2\\
25488&2&$44388/1241$&1\\
25650&2&$2727/73$&1\\
25812&2&$48330/1241$&2\\
25974&2&$50301/1241$&2\\
26298&2&$54243/1241$&2\\
29700&4&$95634/1241$&1\\
30348&4&$103518/1241$&1\\
33264&6&$138996/1241$&1\\\hline
\end{tabular}
\end{table}

\begin{proof}[Construction and exhaustive finite verification]
The standard integral-coordinate Leech model was reconstructed from
the twelve binary Golay generators, from the code construction. Every one of the 654 stabilizer generators was
then checked to preserve its Gram matrix and to have determinant
$\pm1$.

For 123 representatives the intrinsic roots, in their marked
unimodular lattice $N$, generate a rank-two subgroup of
$(\tfrac13N)/N$. Their integral annihilator
\[
 \{x\in N:(x,r)\in\mathbb Z\text{ for every root }r\}
\]
has index nine and is the required perpendicular lattice. Both marked
models were independently built: the $A_2^{12}$ model by the explicit
neighbour construction below, and the $A_1^{24}$ model by binary Golay
glue on 24 orthogonal roots of norm two. The squared norm of each
root was checked in the reconstructed metric.

The six remaining representatives were treated separately. If a matrix
preserves a plane over $\mathbb F_3$, its restriction to that plane
satisfies $(X^8-1)^2$. Intersecting the kernels of this polynomial
for the generators and deterministic words in them gives
spaces of dimension two or three in these six cases. Enumerating all
planes in the latter spaces and testing isotropy and invariance
recovers a unique plane for representatives 2, 6, 127, 128 and 129. For representative
4 there are two candidates; both give $(k,t)=(23058,0)$. Independently,
the defining intersection construction below gives the same pair.
Thus this ambiguity does not affect the theta table. In particular,
the reconstruction does not silently choose one of the two planes.

Each resulting Gram matrix $G$ was checked to have determinant 81,
even diagonal, and positive principal minors. Also $3G^{-1}$ was
checked to be integral with even diagonal. This proves the asserted
level; it cannot be one for a nonunimodular even lattice. The Gram
matrix of $3J^\vee$ is $9G^{-1}$.

For enumeration, an exact unimodular basis change is verified first.
Writing $D_i$ for the principal minors, $D_0=1$, fraction-free
elimination gives integers $m_{ij}$ with
\[
 x^tGx=\sum_{i=0}^{23}
 \frac{(D_{i+1}x_i+\sum_{j>i}m_{ij}x_j)^2}{D_iD_{i+1}}.
\]
The complete matrix identity is checked, not just its diagonal.
Descending in $i$, the residual recursion is
\[
 R_i=\frac{D_iR_{i+1}-(D_{i+1}x_i+
                 \sum_{j>i}m_{ij}x_j)^2}{D_{i+1}}.
\]
An integer square root gives the closed interval of possible $x_i$;
every integer in that interval is visited, and each division is
checked to be exact. Induction on $i$ proves that the enumeration
includes all boundary vectors and omits no vector of the stated norm.
Norm-four counts in $J$ use a new C++ implementation; norm-six counts
in $3J^\vee$ use a separately written Python-integer recursion.
The certificates record the original and reduced Gram matrices,
basis changes, principal minors, elimination coefficients, shell
counts, and rational theta coordinates. They reproduce all entries
of the table. Theorem~\ref{thm:universal-theta-coordinates} then
identifies the entire theta series, rather than merely a truncated
expansion.
\end{proof}

\paragraph{Scope of the classification check.}
The computation above verifies the lattices and theta coordinates of
the representatives. It is not a new proof that the representatives
form all 129 orbits under $\mathrm{Co}_0$, nor a new proof of their
full stabilizer orders. Using the orders, the three exact
incidence checks are
\[
\begin{array}{rcl}
 \sum_P |\mathrm{Co}_0|/|\mathrm{Co}_{0,P}|
   &=&61549107892647466000,\\
 \sum_P (|\mathrm{Co}_0|/|\mathrm{Co}_{0,P}|)k(P)
   &=&196560(3^{22}-1)(3^{20}-1)/16,\\
 \sum_P (|\mathrm{Co}_0|/|\mathrm{Co}_{0,P}|)t(P)
   &=&16773120(3^{10}+1)(3^{11}-1)/2.
\end{array}
\]
The first value is the number of totally isotropic planes in the
split 24-dimensional quadratic space over $\mathbb F_3$. These
equalities corroborate the classification and the shell counts; they do
not replace orbit-completeness or stabilizer proofs.

\section{The defining intersection and stability of the residue action}
We now make explicit which lattices carry the action in
Theorem~\ref{thm:residue-eta}. This stability condition is necessary:
equal rational representations do not themselves imply that a given
integral lattice is preserved.

Use the Euclidean model $A_5^4\perp D_4$ in four sum-zero six-spaces
and one four-space, with
$\lambda=(5,-1,-1,-1,-1,-1)/6$ and
$\omega_j=\sum_{i=1}^j e_i-(j/6)\sum_{i=1}^6e_i$.
The Niemeier lattice $N$ is obtained by adjoining the glue vectors
\[
 (3,3,3,3;0),\ (3,3,0,0;(1,0,0,0)),\
 (3,0,3,0;(\tfrac12,\tfrac12,\tfrac12,\tfrac12)),\
 (2,2,2,0;0),\ (2,4,0,2;0),
\]
where the first four entries multiply $\lambda$. Put
\[
 \rho=((5,3,1,-1,-3,-5)/2)^4\oplus(3,2,1,0),\qquad
 \alpha=(e_1-e_2,0,0,0;0),
\]
and, for $d\mid6$, define
\[
 K_d=\{x\in N:(\rho,x)\in d\mathbb Z\},\qquad
 N_d=K_d+\mathbb Z\bigl(\rho/d-(6/d)\alpha\bigr).
\]
Let $\Lambda_0=N_6$ and $J=\Lambda_0\cap c^2\Lambda_0$.

\begin{theorem}\label{thm:residue-stability}
The full six-cycle $c$ preserves $J$. Among its eight even unimodular
extensions, two have root system $A_2^{12}$ and are individually
preserved by $c$. The other six are Leech lattices, and $c$ permutes
them in two cycles of length three. Consequently:
\begin{enumerate}
\item The full group $S=\langle289\rangle$ acts on $J$ and on both
$A_2^{12}$ extensions; its Frame shapes and eta products are exactly
those in Theorem~\ref{thm:residue-eta}.
\item The subgroup preserving any one of the six Leech extensions
is $\langle289^3\rangle=\{1,169\}$. On such a fixed extension its
two eta products are $\Delta$ and $\eta(2\tau)^{12}$.
\end{enumerate}
For the defining intersection,
\begin{equation}\label{eq:theta-J-exact}
 \Theta_J=\frac{7381}{66430}E(\tau)
 +\frac{59049}{66430}E(3\tau)
 -\frac{264069}{11747}\Delta(\tau)
 -\frac{103276701}{11747}\Delta(3\tau)
 +\frac{14823}{1241}F(\tau).
\end{equation}
Every one of its eight unimodular extensions has $z=0$ in the same
five-dimensional basis. Thus the $F$ coordinate of
$\Theta_M-\Theta_J$ is $-14823/1241$ for every such extension $M$.
\end{theorem}

\begin{proof}
All lattice spans can be computed as integer Hermite forms after
multiplication of ambient coordinates by twelve. The five glue
classes generate an isotropic subgroup of order 72; hence the
determinant of $N$ is $6^4\cdot4/72^2=1$. Its evenness follows
from the integral pairings and half norms of the displayed vectors.
Since $(\rho,\alpha)=1$ and $\rho^2=84$, the map defining $K_d$
is onto. Its index is $d$, and the adjoining vector has norm
$4(6/d)^2$, pairs integrally with $K_d$, and has exact order $d$
modulo $K_d$. Thus $N_d$ is even unimodular.

The implementation forms intersections by the identity
$A\cap B=(A^\vee+B^\vee)^\vee$, using exact inverses and Hermite
forms. It verifies determinant 81 for $J$, equality of the three
pairwise intersections of $\Lambda_0,c^2\Lambda_0,c^4\Lambda_0$,
and invariance under $c^2$. There is also a short direct check for
$c^3$. Let $\beta$ be four copies of $\omega_3$, with zero $D_4$
component. Then $\beta\in N$ by its glue class $(3,3,3,3;0)$, and
$c^3\rho-\rho=-6\beta$. It follows that $c^3K_6=K_6$.
Writing $v=\rho/6-\alpha$, we have
$c^3v-v=-\beta-c^3\alpha+\alpha\in K_6$:
its pairing with $\rho$ is $-18-1+1=-18$.
Thus $c^3\Lambda_0=\Lambda_0$, and, as $c$ commutes with $c^2$,
also $c^3J=J$. Together with $c^2J=J$ this gives $cJ=J$.
The certificate additionally supplies the integral matrix of $c$
on a basis of $J$ and checks the full isometry identity.

Set $u=(\omega_2,\omega_4,0,\omega_2;0)$ and
$(e_1,e_2,e_3,e_4)=(u,v,c^2u,c^2v)$. Exact reduction shows that
their classes form a basis of $J^\vee/J\simeq\mathbb F_3^4$,
with Gram matrix and quadratic form
\[
 \begin{pmatrix}
 4&2&-2&-4/3\\2&4&-2/3&-2/3\\
 -2&-2/3&4&2\\-4/3&-2/3&2&4
 \end{pmatrix},\qquad
 Q(a,b,c,d)=\frac{2ad+bc+bd}{3}\pmod{\mathbb Z}.
\]
The matrix of $c$ on row coordinates in this quotient is
\[
 A=\begin{pmatrix}
 2&0&2&0\\1&2&2&2\\1&0&0&0\\0&1&0&0
 \end{pmatrix}\pmod3.
\]
In particular $A^3=1$. Even unimodular extensions are in bijection
with maximal totally isotropic planes of this split four-space.
Enumeration of its 81 vectors gives eight such planes. In the
certificate ordering, their root counts and $c$ permutation are
\[
 (0,0,72,0,0,72,0,0),\qquad
 (0\ 3\ 4)(1\ 7\ 6),
\]
with 2 and 5 fixed. The root counts come from exhaustive norm-two
enumeration. The rank-24 even unimodular classification identifies
the six rootless lattices as Leech lattices and the two with 72 roots
as $A_2^{12}$. This proves all the stability assertions.

Finally, independent enumeration for $J$ gives
\[
 \Theta_J=1+23058q^2+1849032q^3+44291124q^4+O(q^5),
 \qquad r_{3J^\vee}(6)=0.
\]
The norm-eight enumeration visited 261581721 recursion nodes and
used no floating-point inclusion decisions. Substitution into
Theorem~\ref{thm:universal-theta-coordinates} gives
\eqref{eq:theta-J-exact}; the additional $q^3,q^4$ counts also
verify it independently by the Sturm bound four.
For any even unimodular rank-24 lattice $M$, its level-one theta
series is $E+(r_M(2)-C)\Delta$. The six Leech cases therefore have
coordinates $(1,0,-C,0,0)$, and the two rooted cases have
$(1,0,72-C,0,0)$. The statements about $z$ follow.
\end{proof}

\section{Shell Dirichlet series and what Dirichlet GRH controls}
For $\Theta_J=1+\sum_{n\ge1}r_J(2n)q^n$, put
$D_J(s)=\sum_{n\ge1}r_J(2n)n^{-s}$. All formulas in this section
first hold for $\Re s>12$ and then continue meromorphically.
The decomposition above gives
\begin{equation}\label{eq:J-shell-series}
 D_J(s)=C(a+b3^{-s})\zeta(s)\zeta(s-11)
       +(x+y3^{-s})L(\Delta,s)+zL(F,s).
\end{equation}
Let $\chi$ be the nontrivial character modulo three and let $\chi_0$
be the principal character modulo three, both extended by zero on
multiples of three. Coefficientwise twisting gives
\[
 D_{J,\chi}(s)=aC L(s,\chi)L(s-11,\chi)
                  +xL(\Delta\otimes\chi,s)
                  +zL(F\otimes\chi,s),
\]
where the two terms supported on multiples of three disappear.
For the principal twist,
\begin{align*}
 D_{J,\chi_0}(s)={}&aC(1-3^{-s})(1-3^{11-s})
                        \zeta(s)\zeta(s-11)\\
 &+x(1-252\,3^{-s}+3^{11-2s})L(\Delta,s)
   +z(1+243\,3^{-s})L(F,s).
\end{align*}
These finite factors use $\tau(3)=252$ and $a_3(F)=-243$.

Dirichlet GRH locates the nontrivial zeros of the Dirichlet
$L$ factors, so the nontrivial zeros contributed by those factors
lie on $\Re s=1/2$ or $23/2$. It says nothing by itself about
the cuspidal factors or the zeros of the sums above. There are also
elementary zeros that must not be discarded. In particular the
untwisted Eisenstein summand has zeros at
\[
 s=\log_3\!\left(\frac{59049}{7381}\right)
       +\frac{(2j+1)\pi i}{\log3},\qquad j\in\mathbb Z,
\]
coming from $a+b3^{-s}$. The principal twist has additional finite-factor
zeros on $\Re s=0$ and $\Re s=11$. These assertions concern the
specified summands: adding cusp-form $L$ series can move or remove
such zeros. No transfer of a Riemann-hypothesis statement to
$D_J$ or its twists follows from linear decomposition alone.

\section{All weight-one Frame-shape forms}
The weight-one list separates a small cuspidal family from a larger
Eisenstein family, and requires the factors at primes dividing the
level to be displayed.
For a Frame shape \(\pi=\prod_d d^{r_d}\), define
\(\eta_\pi(\tau)=\prod_d\eta(d\tau)^{r_d}\). Its weight is
\(\tfrac12\sum_d r_d\).
Kond\=o's complete table of the Leech automorphism group gives
29 class entries of weight one, with 28 distinct Frame shapes
\cite[Table~I, pp.~355--357]{Kondo1985}.
The signs on class names below follow his convention on p.~354.

Let \(d_1,d_2\) be fundamental discriminants, allowing \(1\) for
the trivial character. Assume their characters have opposite parity:
\(\chi_{d_1}(-1)\chi_{d_2}(-1)=-1\). Write
\[
 E_{d_1,d_2}(\tau)=c_0+
 \sum_{n\ge1}\left(\sum_{a\mid n}
             \chi_{d_1}(a)\chi_{d_2}(n/a)\right)q^n.
\]
If neither character is trivial, \(c_0=0\); if one is trivial
and the other is \(\chi_d\), put \(c_0=L(0,\chi_d)/2\).
These are the weight-one Eisenstein series, with
\[
 L(E_{d_1,d_2},s)=L(s,\chi_{d_1})L(s,\chi_{d_2}),
\]
where the constant term is omitted in the Dirichlet transform.

\begin{theorem}\label{thm:weight-one-list}
The distinct weight-one eta products for \(M_{24}\) are
\[
 \eta(\tau)\eta(23\tau),\qquad
 \eta(3\tau)\eta(21\tau),\qquad
 \eta(12\tau)^2.
\]
They are irreducible dihedral cusp forms.
For \(\mathrm{Co}_0\), the five distinct cuspidal products are
those in Table~\ref{tab:weight-one-cusp}; all other 23 distinct
products are given in the following table.
For each noncuspidal row, if \(P(s)=\sum_t c_t t^{-s}\), then
\begin{equation}\label{eq:weight-one-eis}
 \eta_\pi(\tau)=\sum_t c_tE_{d_1,d_2}(t\tau),\qquad
 L(\eta_\pi,s)=P(s)L(s,\chi_{d_1})L(s,\chi_{d_2}).
\end{equation}
The primitive Galois representation in that row is
\(\chi_{d_1}\oplus\chi_{d_2}\), hence reducible.
Its Hecke character over the imaginary quadratic field of the
listed discriminant \(D\) is
\(\chi_{d_1}\circ\mathrm N_{K/\Q}\), equivalently
\(\chi_{d_2}\circ\mathrm N_{K/\Q}\).
The factors \(P\) are part of the displayed eta-product
\(L\)-series and cannot be discarded.
\end{theorem}

\begin{table}[htbp]
\centering
\caption{The five cuspidal weight-one products for \(\mathrm{Co}_0\).
The chosen CM realization is \(K=\Q(\sqrt{D_K})\), with
Hecke conductor \(f\mathcal O_K\).
The character is a faithful character of the indicated cyclic ring class
group, up to inversion.}
\label{tab:weight-one-cusp}
\begin{tabular}{@{}lrrrr@{}}
\toprule
Eta product&Level&\(D_K\)&\(f\)&Character order\\\midrule
\(\eta(\tau)\eta(23\tau)\)&23&\(-23\)&1&3\\
\(\eta(2\tau)\eta(22\tau)\)&44&\(-11\)&2&3\\
\(\eta(3\tau)\eta(21\tau)\)&63&\(-7\)&3&4\\
\(\eta(4\tau)\eta(20\tau)\)&80&\(-20\)&2&4\\
\(\eta(12\tau)^2\)&144&\(-4\)&6&4\\
\bottomrule
\end{tabular}
\end{table}

{\small
\begin{longtable}{@{}llrcll@{}}
\caption{The noncuspidal weight-one Frame shapes. Here \(X_p=p^{-s}\).}\\
\toprule Class&Frame shape&Level&\(\chi\)&\((d_1,d_2)\)&\(P(s)\)\\\midrule\endfirsthead
\toprule Class&Frame shape&Level&\(\chi\)&\((d_1,d_2)\)&\(P(s)\)\\\midrule\endhead
-6F & $\frac{1\cdot6^{6}}{2^{2}\cdot3^{3}}$ & 6 & $\chi_{-3}$ & $(1,-3)$ & $1 - X_{2}$ \\
8D & $\frac{8^{4}}{4^{2}}$ & 8 & $\chi_{-4}$ & $(1,-4)$ & $1 - X_{2}$ \\
+9B & $\frac{9^{3}}{3}$ & 9 & $\chi_{-3}$ & $(1,-3)$ & $1 - X_{3}$ \\
-12D & $\frac{2\cdot3^{3}\cdot12^{3}}{1\cdot4\cdot6^{3}}$ & 12 & $\chi_{-3}$ & $(1,-3)$ & $- (X_{2} - 1) (2 X_{2} + 1)$ \\
12G & $\frac{4^{2}\cdot12^{2}}{2\cdot6}$ & 12 & $\chi_{-3}$ & $(1,-3)$ & $- (X_{2} - 1) (X_{2} + 1)$ \\
+12H & $\frac{2^{3}\cdot6\cdot12^{2}}{1\cdot3\cdot4^{2}}$ & 12 & $\chi_{-4}$ & $(1,-4)$ & $1 - X_{3}$ \\
-12K & $\frac{1^{3}\cdot12^{3}}{2\cdot3\cdot4\cdot6}$ & 12 & $\chi_{-3}$ & $(1,-3)$ & $(X_{2} - 1) (2 X_{2} - 1)$ \\
-14B & $\frac{2^{2}\cdot14^{2}}{1\cdot7}$ & 14 & $\chi_{-7}$ & $(1,-7)$ & $1 - X_{2}$ \\
+15E & $\frac{1^{2}\cdot15^{2}}{3\cdot5}$ & 15 & $\chi_{-15}$ & $(-3,5)$ & $1$ \\
16A & $\frac{2^{2}\cdot16^{2}}{4\cdot8}$ & 32 & $\chi_{-8}$ & $(-4,8)$ & $1$ \\
-16B & $\frac{1^{2}\cdot16^{2}}{2\cdot8}$ & 16 & $\chi_{-4}$ & $(1,-4)$ & $(X_{2} - 1) (2 X_{2} - 1)$ \\
-18B & $\frac{1^{2}\cdot9\cdot18}{2\cdot3}$ & 18 & $\chi_{-3}$ & $(1,-3)$ & $(2 X_{2} - 1) (X_{3} - 1)$ \\
+18C & $\frac{1\cdot2\cdot18^{2}}{6\cdot9}$ & 18 & $\chi_{-3}$ & $(1,-3)$ & $(X_{2} - 1) (3 X_{3} - 1)$ \\
-18C & $\frac{2^{2}\cdot9\cdot18}{1\cdot6}$ & 18 & $\chi_{-3}$ & $(1,-3)$ & $- (X_{2} + 1) (X_{3} - 1)$ \\
+20C & $\frac{1\cdot2\cdot10\cdot20}{4\cdot5}$ & 20 & $\chi_{-20}$ & $(-4,5)$ & $1$ \\
-20C & $\frac{2^{2}\cdot5\cdot20}{1\cdot4}$ & 20 & $\chi_{-4}$ & $(1,-4)$ & $1 - X_{5}$ \\
24E & $\frac{2\cdot6\cdot8\cdot24}{4\cdot12}$ & 48 & $\chi_{-3}$ & $(-4,12)$ & $1$ \\
+24F & $\frac{1\cdot4\cdot6\cdot24}{3\cdot8}$ & 24 & $\chi_{-24}$ & $(-3,8)$ & $1$ \\
-24F & $\frac{2\cdot3\cdot4\cdot24}{1\cdot8}$ & 24 & $\chi_{-8}$ & $(1,-8)$ & $1 - X_{3}$ \\
-28A & $\frac{1\cdot4\cdot7\cdot28}{2\cdot14}$ & 28 & $\chi_{-7}$ & $(1,-7)$ & $(X_{2} - 1) (2 X_{2} - 1)$ \\
+30D & $\frac{1\cdot6\cdot10\cdot15}{3\cdot5}$ & 30 & $\chi_{-15}$ & $(-3,5)$ & $X_{2} + 1$ \\
-30D & $\frac{2\cdot3\cdot5\cdot30}{1\cdot15}$ & 30 & $\chi_{-15}$ & $(1,-15)$ & $1 - X_{2}$ \\
-30E & $\frac{2\cdot3\cdot5\cdot30}{6\cdot10}$ & 30 & $\chi_{-15}$ & $(-3,5)$ & $2 X_{2} + 1$ \\
\bottomrule\end{longtable}

}

\begin{proof}
For completeness, the finite classification and the identities can
be checked by short exact computations, with no identification from
approximate prime coefficients.
Selecting \(\sum r_d=2\) from the three pages of Kond\=o's
table gives exactly the rows displayed here. The only duplicate
Frame shape is \(1\cdot23\), for \(+23A,+23B\).
The \(M_{24}\) cycle-shape table selects the three products in
the statement \cite[Table~I, p.~154]{Koike1985}.

For every row, \(\sum d r_d=24\), so the eta quotient begins
with \(q\). With the level \(N\) in the table, its second eta
congruence is
\(\sum (N/d)r_d\equiv0\pmod{24}\).
Its character is the Kronecker character of the square class
\(-\prod d^{r_d}\), giving exactly the displayed \(\chi_D\).
For a cusp with denominator \(c\mid N\), the order is
\begin{equation}\label{eq:eta-cusp-order-weight1}
 \frac{N}{24\,c\,\gcd(c,N/c)}
       \sum_d\frac{\gcd(c,d)^2}{d}r_d.
\end{equation}
Substitution gives nonnegative orders for every row; all orders
are positive exactly for the five products in
Table~\ref{tab:weight-one-cusp}.
Thus the eta quotients and their proposed Eisenstein expressions
are holomorphic forms of the same weight, level and character.

We recall why a finite coefficient comparison proves these identities.
For a holomorphic form of weight \(k\) on \(\Gamma_0(N)\), the
valence formula applied to the product of its translates to
\(\mathrm{SL}_2(\Z)\) bounds its order at infinity by
\(k[\mathrm{SL}_2(\Z):\Gamma_0(N)]/12\), unless it is zero.
A power of that product removes its finite character if necessary.
Consequently two such forms agree if their coefficients through
\[
 B=\left\lfloor\frac{kN}{12}\prod_{p\mid N}(1+1/p)\right\rfloor
\]
agree. For the noncuspidal rows, the coefficients in
Table~\ref{tab:weight-one-checks} follow on the eta side by
multiplying \(\prod_{d,n}(1-q^{dn})^{r_d}\), and on the
Eisenstein side by the divisor sum in its definition.
The constant terms vanish: either both characters are nontrivial,
or \(\sum c_t=0\). This proves \eqref{eq:weight-one-eis}.
It also makes the character and level information directly
checkable.

Here is a separate exact identification of the five cusp forms.
For \(Q=[a,b,c]\), put
\(\theta_Q=\sum_{x,y\in\Z}q^{ax^2+bxy+cy^2}\).
The following identities hold:
\[
\begin{array}{c|cc|r}
 N&Q_0&Q_1&B\\\hline
23 &[1,1,6]&[2,1,3]&2\\
44 &[1,0,11]&[3,2,4]&6\\
63 &[1,1,16]&[4,1,4]&8\\
80 &[1,0,20]&[4,0,5]&12\\
144&[1,0,36]&[4,0,9]&24
\end{array}
\qquad
 \eta_\pi=\frac{\theta_{Q_0}-\theta_{Q_1}}2.
\]
Binary theta modularity follows from Poisson summation, with
level \(N\) and character \(\chi_{-N}\). The two sides have
the following common initial polynomials, with terms beyond
the indicated \(B\) omitted:
\[
\begin{array}{c|l}
23&q-q^2\\
44&q-q^3-q^5\\
63&q-q^4-q^7\\
80&q-q^5-q^9\\
144&q-2q^{13}
\end{array}
\]
These are obtained by enumerating the pairs \((x,y)\) in the
two positive quadratic forms; the same valence bound proves
all five identities.

Reduction of primitive positive binary quadratic forms gives
the complete class lists
\[
\begin{array}{c|l}
-23 &[1,1,6],\ [2,\pm1,3]\\
-44 &[1,0,11],\ [3,\pm2,4]\\
-63 &[1,1,16],\ [2,\pm1,8],\ [4,1,4]\\
-80 &[1,0,20],\ [3,\pm2,7],\ [4,0,5]\\
-144&[1,0,36],\ [5,\pm4,8],\ [4,0,9].
\end{array}
\]
The first two class groups have order three. In each of the
last three, only the principal class and the final class are
self-inverse, so the group of order four is cyclic.
For a faithful character \(\xi\) of one of these groups,
counting ideals of each class gives
\(\theta_\xi=\tfrac12\sum_{[Q]}\xi([Q])\theta_Q\).
Inverse classes have the same theta series. In order three their
two character values sum to \(-1\); in order four their values
cancel, and the order-two class has value \(-1\).
Thus \(\theta_\xi\) is precisely the displayed theta difference.
The characters are primitive at the conductors in the table:
the class groups of every proper divisor conductor have order
smaller than the stated character order. Since conjugation
sends \(\xi\) to \(\xi^{-1}\ne\xi\), their induced
two-dimensional representations are irreducible.
This proves the cuspidal and Hecke-character assertions.

For an Eisenstein row, let \(\chi_K=\chi_D\).
The table gives \(\chi_{d_2}=\chi_{d_1}\chi_K\) as primitive
characters away from their conductors. Restricting to \(G_K\)
identifies the two characters. Quadratic induction therefore
decomposes as
\[
 \operatorname{Ind}_{G_K}^{G_\Q}
        (\chi_{d_1}|_{G_K})
   =\chi_{d_1}\oplus\chi_{d_2}.
\]
The ideal-theoretic character is
\(\chi_{d_1}\circ\mathrm N_{K/\Q}\).
Its conductor norm is
\(|d_1d_2|/|D|\), which is \(1\) except for the row \(16A\)
(norm \(4\)) and \(24E\) (norm \(16\)).
The corresponding conductor ideals are \((2)\) in
\(\Q(\sqrt{-2})\), and \((4)\) in \(\Q(\sqrt{-3})\).
This proves the final assertion, including the ramified cases.
\end{proof}

{\small
\begingroup
\renewcommand{\tablename}{Table}
\begin{longtable}{@{}lrl@{}}
\caption{Exact coefficients through the weight-one Sturm bound \(B\).
The constant coefficient is zero in every row.}\label{tab:weight-one-checks}\\
\toprule Class&\(B\)&\((a_1,\ldots,a_B)\)\\\midrule\endfirsthead
\toprule Class&\(B\)&\((a_1,\ldots,a_B)\)\\\midrule\endhead
-6F & 1 & $(1)$ \\
8D & 1 & $(1)$ \\
+9B & 1 & $(1)$ \\
-12D & 2 & $(1,1)$ \\
12G & 2 & $(1,0)$ \\
+12H & 2 & $(1,1)$ \\
-12K & 2 & $(1,-3)$ \\
-14B & 2 & $(1,1)$ \\
+15E & 2 & $(1,-2)$ \\
16A & 4 & $(1,0,-2,0)$ \\
-16B & 2 & $(1,-2)$ \\
-18B & 3 & $(1,-2,0)$ \\
+18C & 3 & $(1,-1,-2)$ \\
-18C & 3 & $(1,1,0)$ \\
+20C & 3 & $(1,-1,-2)$ \\
-20C & 3 & $(1,1,0)$ \\
24E & 8 & $(1,0,-1,0,0,0,-2,0)$ \\
+24F & 4 & $(1,-1,-1,1)$ \\
-24F & 4 & $(1,1,1,1)$ \\
-28A & 4 & $(1,-1,0,-1)$ \\
+30D & 6 & $(1,-1,-1,1,-1,1)$ \\
-30D & 6 & $(1,1,1,1,1,1)$ \\
-30E & 6 & $(1,0,-1,-1,-1,0)$ \\
\bottomrule\end{longtable}

\endgroup
}

The qualification at bad primes is visible already in the row \(-12K\).
Its coefficients satisfy \(a(2)=-3\) and \(a(4)=3\), so it is
not an eigenfunction of \(U_2\): the required relation would be
\(a(4)=a(2)^2=9\). It is an Eisenstein eigenfunction at every
prime not dividing the level. Thus the assertion that every
\(\mathrm{Co}_0\) Frame-shape eta quotient is a simultaneous
newform is false.

For the reducible rows, Alp\"oge's theorem puts all nontrivial
zeros of the two primitive Dirichlet factors on \(\Re s=1/2\).
The factors \(P\) in the table also have zeros, on \(\Re s=0\)
or \(\Re s=1\), with possible cancellation at a pole of a
Dirichlet factor. For example, \(-12K\) has
\[
 P(s)=(1-2^{-s})(1-2^{1-s}),
\]
so it has infinitely many elementary zeros on \(\Re s=1\);
the zero at \(s=1\) itself cancels the zeta pole.
The five primitive dihedral cusp-form \(L\)-functions are
Hecke \(L\)-functions over quadratic fields, not products of
Dirichlet \(L\)-functions. Their Riemann hypotheses do not
follow from the Dirichlet theorem.

\section{The completed mock forms and their Mellin equations}
\label{sec:mock-mellin}
The completed mock forms have a modular functional equation, but their
holomorphic parts need not have an ordinarily convergent Dirichlet series.
The appropriate general framework uses test functions, as in Diamantis,
Lee, Raji and Rolen \cite[Definition~4.2, Lemma~4.4 and Theorem~4.5]{DLRR2023}.
We give both that equation and a specified regularized Mellin version.

Let $\widehat H$ be an umbral completion of weight $1/2$, with the odd
component convention $H_{-r}=-H_r$.  Write $G=\xi_{1/2}\widehat H$, where
\[
 \xi_kF=2iy^k\overline{\frac{\partial F}{\partial\bar\tau}}.
\]
This defines the normalization of $G$; it is a fixed scalar multiple of the
unary shadow $S^X$.  With this convention the nonholomorphic part is
\begin{equation}\label{eq:mock-nonholomorphic-part}
 H^-(\tau)=-\sum_{\nu>0}\overline{b(\nu)}(4\pi\nu)^{-1/2}
       \Gamma(1/2,4\pi\nu y)q^{-\nu},
 \qquad G(\tau)=\sum_{\nu>0}b(\nu)q^\nu.
\end{equation}
The values of $\nu$ may be rational.  The normalization is useful because
applying $\xi_{1/2}$ to each summand gives exactly $b(\nu)q^\nu$.

For the untwined form, put $N=2m$ and define, on the odd subspace,
\begin{equation}\label{eq:mock-weil-involution}
 W_{r,r'}=\frac{i}{\sqrt N}e^{2\pi irr'/N}.
\end{equation}
The Jacobi transformation law and theta decomposition give
\begin{equation}\label{eq:mock-axis-transformation}
 \widehat H(i/y)=y^{1/2}W\widehat H(iy),\qquad W^2=I.
\end{equation}
Indeed the Fourier matrix squared is $N$ times $r\mapsto-r$, which is
$-I$ on odd vectors; the factor $i^2=-1$ gives $W^2=I$.
At $m=6$, equation~\eqref{eq:mock-weil-involution} is the explicit
$12\times12$ formula $W_{r,r'}=i e^{2\pi irr'/12}/\sqrt{12}$,
restricted to $H_0=H_6=0$ and $H_{12-r}=-H_r$.
The same argument applies at every lambency.  For a twining whose modular
group has several relevant cusps, one uses the vector of cusp expansions
and the corresponding matrix in place of $W$.

\begin{proposition}[Test-function equation]\label{prop:mock-test-function}
For a smooth compactly supported function $\phi:(0,\infty)\to\mathbb C$, put
\[
 \mathcal L_{\widehat H}(\phi)=
       \int_0^\infty\widehat H(iy)\phi(y)\,dy,
 \qquad \phi^\vee(y)=y^{-3/2}\phi(1/y).
\]
Then
\begin{equation}\label{eq:mock-test-functional}
 \mathcal L_{\widehat H}(\phi)
       =W\mathcal L_{\widehat H}(\phi^\vee).
\end{equation}
\end{proposition}
\begin{proof}
All integrals converge because the support is a compact subset of
$(0,\infty)$.  Substitute $y=1/t$ and use
\eqref{eq:mock-axis-transformation}.  The Jacobian $t^{-2}$ combines with
$t^{1/2}$ to give $t^{-3/2}$, which is precisely the displayed dual test
function.
\end{proof}

A plain power $\phi(y)=y^{s-3/2}$ is not an admissible test function when
$\widehat H$ has a polar Fourier term: that term grows exponentially as
$y\to\infty$.  To make the Mellin notation unambiguous, fix the following
regularization.  Write the finitely many growing and constant terms at
infinity as
$P(y)=\sum_{\nu<0}a(\nu)e^{-2\pi\nu y}+a(0)$ and put
\begin{align}
 F_H(u)={}&\int_1^\infty\bigl(\widehat H(iy)-P(y)\bigr)y^{u-1}\,dy
 -\frac{a(0)}u\notag\\
 &+\sum_{\nu<0}a(\nu)(2\pi\nu+i0)^{-u}
                 \Gamma(u,2\pi\nu+i0).\label{eq:mock-tail-regularization}
\end{align}
Here $+i0$ specifies continuation from the upper half-plane and thereby
fixes the branch on the negative real axis.  The integral is entire in
$u$ because its integrand decays exponentially at infinity; the other
terms are meromorphic.  Define the shifted, symmetrically regularized
Mellin transform by
\begin{equation}\label{eq:mock-regularized-Mellin}
 \mathcal M_{\widehat H}(s)=F_H(s-1/2)+WF_H(1-s).
\end{equation}
This is the regularized form of
$\int_0^\infty\widehat H(iy)y^{s-3/2}\,dy$ obtained by splitting at $1$
and using the modular transformation on the lower interval.
Its functional equation is now precise:
\begin{equation}\label{eq:mock-Mellin-FE}
 \mathcal M_{\widehat H}(s)
       =W\mathcal M_{\widehat H}(3/2-s).
\end{equation}
It follows immediately from $W^2=I$ and
\eqref{eq:mock-regularized-Mellin}.

\begin{proposition}[The shadow contribution]\label{prop:mock-shadow-Mellin}
Let $\overline G(\tau)=\sum_{\nu>0}\overline{b(\nu)}q^\nu$ and set
\begin{equation}\label{eq:mock-kernel}
 C(s)=(4\pi)^{-s}\frac{\Gamma(s)}{s-1/2}
 {}_2F_1\left(s-\tfrac12,s;s+\tfrac12;\tfrac12\right).
\end{equation}
For $\Re s>1$, the ordinary Mellin transform of $H^-$ converges absolutely
and is
\begin{equation}\label{eq:mock-shadow-Mellin}
 \int_0^\infty H^-(iy)y^{s-3/2}\,dy=-C(s)L(\overline G,s).
\end{equation}
Let $\mathcal M_{H^+}(s)$ denote the holomorphic-part Mellin transform with
the regularization fixed by
\eqref{eq:mock-tail-regularization}--\eqref{eq:mock-regularized-Mellin} and
linearity.  Then, by meromorphic continuation,
\begin{align}
 \mathcal M_{H^+}(s)-C(s)L(\overline G,s)
 =W\bigl\{\mathcal M_{H^+}(3/2-s)
       -C(3/2-s)L(\overline G,3/2-s)\bigr\}.
 \label{eq:mock-inhomogeneous-Mellin}
\end{align}
Every shadow $L$-function here is explicitly a sum of the
$L(2s-1,\psi)$ in Theorem~\ref{thm:all-shadow-characters}.
\end{proposition}
\begin{proof}
For a single term of \eqref{eq:mock-nonholomorphic-part}, set $t=4\pi\nu y$.
The resulting scalar integral, without the coefficient
$-\overline{b(\nu)}(4\pi\nu)^{-s}$, is
\[
 K(s)=\int_0^\infty t^{s-3/2}e^{t/2}\Gamma(1/2,t)\,dt.
\]
For $\Re s>1/2$, insert
$\Gamma(1/2,t)=\int_t^\infty e^{-v}v^{-1/2}\,dv$, expand $e^{t/2}$, and
integrate first over $0<t<v$.  Absolute convergence gives
\begin{align*}
 K(s)&=\sum_{j\geq0}\frac{\Gamma(s+j)}{2^j j!(s+j-1/2)}\\
 &=\frac{\Gamma(s)}{s-1/2}
 {}_2F_1\left(s-\tfrac12,s;s+\tfrac12;\tfrac12\right).
\end{align*}
For a unary shadow, $\sum_{\nu>0}|b(\nu)|\nu^{-\sigma}$ converges for
$\sigma>1$, so summation and integration are justified in the stated
half-plane.  This proves \eqref{eq:mock-shadow-Mellin}.
Subtract this convergent term from the fixed regularized transform of
$\widehat H$ and continue meromorphically.  Substitution into
\eqref{eq:mock-Mellin-FE} proves
\eqref{eq:mock-inhomogeneous-Mellin}.
\end{proof}

The Dirichlet Riemann hypothesis therefore controls specified constituents
of the inhomogeneous shadow term.  It gives no zero-location theorem for
$\mathcal M_{H^+}$ or for $\mathcal M_{\widehat H}$ by this identity.
Moreover, the naive expression $\sum a(\nu)\nu^{-s}$ for a holomorphic
mock part may fail to converge in every half-plane because polar modular
forms can have coefficients of exponential order in $\sqrt\nu$.
A question about its zeros must first specify a regularization or a
test-function family.  The choice in
\eqref{eq:mock-tail-regularization} is explicit; it is not an Euler product,
and equation~\eqref{eq:mock-inhomogeneous-Mellin} provides no argument that
its zeros lie on a vertical line.

\section{Fricke constants as glue orders}
This section explains why the numerical constants in the Fricke involutions
are the indices of the root lattices in their Niemeier extensions.
For an even unimodular lattice \(N\) containing a full-rank root lattice
\(R\), the index formula gives
\[
 [N:R]^2=\det R.
\]
If \(R=A_{n-1}^{e}\), \(e=24/(n-1)\), then \(\det R=n^e\).
Consequently the glue order is
\[
 C_n=n^{12/(n-1)}.
\]
Put \(T_n(\tau)=(\eta(\tau)/\eta(n\tau))^e\).
The root systems and the relevant genus-zero functions are listed in
Cheng--Duncan--Harvey \cite[Sect.~2.2, Table~1]{CDH2014}.

\begin{theorem}\label{thm:fricke-glue}
For \(n\in\{2,3,4,5,7,9,13,25\}\), the normalized Hauptmodul of
the group generated by \(\Gamma_0(n)\) and the Fricke involution \(W_n\) is
\[
 J_n^+=T_n+e+\frac{C_n}{T_n}.
\]
The corresponding Monster classes, in the same order, are
\(2A,3A,4A,5A,7A,9A,13A,25A\).
For the root system \(A_5^4D_4\), put
\[
 T_6=\frac{\eta(\tau)^5\eta(3\tau)}
              {\eta(2\tau)\eta(6\tau)^5}.
\]
The glue order is \(72\), and the normalized Hauptmodul
\[
 J_6=T_6+5+\frac{72}{T_6}
\]
is the McKay--Thompson series \(6B\). Its group is
\(\langle\Gamma_0(6),W_6\rangle\).
\end{theorem}
\begin{proof}
The transformation \(\eta(-1/\tau)=\sqrt{-i\tau}\eta(\tau)\)
gives \(T_n(W_n\tau)=C_n/T_n(\tau)\).
The eta congruences and square character show that \(T_n\) is a
modular function on \(X_0(n)\). At a cusp with denominator \(c\mid n\)
its order is
\[
 \frac{e}{24}\frac{n}{\gcd(c^2,n)}
                   \left(1-\frac{c^2}{n}\right).
\]
For every \(n\) in the statement this is \(1\) at \(c=1\),
\(-1\) at \(c=n\), and zero at any intermediate denominator.
There are no zeros or poles in the upper half-plane, since eta
does not vanish there. Thus the divisor is \((0)-(\infty)\);
a function with a unique simple pole has degree one and identifies
\(X_0(n)\) with the projective line. Its expansion begins
\(q^{-1}-e+O(q)\).
For an indeterminate \(t\), the fixed field of
\(t\mapsto C_n/t\) in \(\C(t)\) is \(\C(t+C_n/t)\):
\(t\) satisfies \(t^2-(t+C_n/t)t+C_n=0\), and the
involution is nontrivial. Thus the displayed invariant generates
the function field of the quotient and has expansion \(q^{-1}+O(q)\).
This proves the normalized Hauptmodul assertion.

For \(A_5^4D_4\), the determinant is \(6^4\cdot4\), with square root
\(72\). Applying the same eta transformation to the four factors
of \(T_6\) gives \(T_6(W_6\tau)=72/T_6(\tau)\).
The cusp orders of \(T_6\), at denominators \(1,2,3,6\), are
\(1,0,0,-1\). The same degree-one argument makes \(T_6\) a
coordinate on \(X_0(6)\).
Here \(T_6=q^{-1}-5+6q+4q^2+O(q^3)\), so the same fixed-field
argument proves the formula for \(J_6\).
The Monster identifications, including the Fricke-only group \(6+6\)
for \(6B\), are those of Conway--Norton
\cite[Table~2, pp.~327--328]{CN1979}.
At level \(25\) one must distinguish the Fricke quotient \(25A\)
from the function for \(\Gamma_0(25)\), traditionally called \(25Z\).
\end{proof}

\begin{center}
\begin{tabular}{c|rrrrrrrr}
\(n\)&2&3&4&5&7&9&13&25\\ \hline
\(e\)&24&12&8&6&4&3&2&1\\
\(C_n\)&4096&729&256&125&49&27&13&5
\end{tabular}
\end{center}

\section{The cusp cover of \(X_0(6)\)}
The four branch values of the cusp cover give a genus-one curve
whose relation with the conductor-\(24\) modular curve can be proved
by rational functions.

\begin{theorem}\label{thm:cusp-isogeny}
The smooth projective completion of
\[
 C:\quad y^2=t(t+8)(t+9),
\]
with its point at infinity as origin, admits a degree-two isogeny to
\[
 E:\quad v^2=(x-1)(x-2)(x+2)=x^3-x^2-4x+4.
\]
On the affine open set \(t\ne-9\), the isogeny is
\begin{equation}\label{eq:cusp-isogeny}
 x=\frac{t^2+16t+72}{4(t+9)},\qquad
 v=\frac{y(t+6)(t+12)}{8(t+9)^2}.
\end{equation}
Its kernel is \(\{O,(-9,0)\}\), and its pullback of \(dx/(2v)\)
is \(2\,dt/(2y)\). The target is Cremona's curve \(24a1\),
whose LMFDB label is \(24.a4\). Thus the Hasse--Weil \(L\)-function
of \(C\) is that of \(\eta(2\tau)\eta(4\tau)\eta(6\tau)\eta(12\tau)\),
including all finite local factors.
\end{theorem}
\begin{proof}
Set \(u=t+9\). Then \(C\) becomes
\[
 y^2=u^3-10u^2+9u.
\]
For a curve \(y^2=u^3+au^2+bu\), direct substitution shows that
\[
 X=u+a+\frac bu,\qquad Y=y\left(1-\frac b{u^2}\right)
\]
satisfy
\[
 Y^2=X^3-2aX^2+(a^2-4b)X.
\]
With \(a=-10,b=9\), this is \(Y^2=X^3+20X^2+64X\).
The change \(X=4(x-2),Y=8v\) gives \(E\) and
the functions in \eqref{eq:cusp-isogeny}.
The map extends over the smooth projective curves. It sends the origin
to the origin, so is a group homomorphism. The rational function \(X\)
has degree two as a function of \(u\); equivalently,
\(u^2+(a-X)u+b=0\), and the kernel is precisely the indicated
two-torsion subgroup. This proves the degree and kernel.
Differentiating the two formulas gives the differential assertion.

For completeness, the dual map, in the same coordinates, is
\[
 t=x-6+\frac4{x-2},\qquad
 y=v\left(1-\frac4{(x-2)^2}\right).
\]
The two compositions are multiplication by two.
The equation and the conductor-\(24\) identification of \(E\)
are recorded in \cite{LMFDB24}. The eta quotient in the statement
has weight two, trivial character and positive orders at every
cusp of \(\Gamma_0(24)\); the two eta congruence sums equal \(24\),
and the product entering the character is \(576\).
The space \(S_2(\Gamma_0(24))\) has dimension one and its old space
is zero. The quotient is therefore the normalized newform of \(E\).
An isogeny identifies the rational Tate modules, and hence their
inertia-invariant Frobenius characteristic polynomials at every prime.
This gives equality of all finite local factors.
\end{proof}

\section{CM zeros of Niemeier theta series}
The zeros at imaginary quadratic arguments are decided by the rational
singular moduli, rather than by a Riemann hypothesis for a Mellin transform.
Write
\(\Theta_L(\tau)=\sum_{v\in L}q^{(v,v)/2}\), let \(N_2(L)\) be the
number of roots, and set \(h=N_2(L)/24\); for the Leech lattice
\(\Lambda\), this convention gives \(h=0\).

\begin{theorem}\label{thm:cm-theta-zeros}
Among the \(24\) positive-definite even unimodular lattices of rank \(24\),
\(\Theta_L\) vanishes at an imaginary quadratic argument only for
the root systems \(E_8^3\) and \(D_{16}E_8\). In both cases the
zero arguments are exactly the \(\mathrm{SL}_2(\Z)\)-orbit of
\(\rho=e^{2\pi i/3}\).
The genus-two theta series of the Leech lattice does not vanish at
any point on the Humbert surface \(H_1\) represented by a pair
of imaginary quadratic arguments.
\end{theorem}
\begin{proof}
The two-dimensional space of weight-twelve modular forms on
\(\mathrm{SL}_2(\Z)\) gives
\begin{equation}\label{eq:niemeier-j}
 \Theta_L=E_4^3+(24h-720)\Delta
          =\Delta\bigl(j-720+24h\bigr),
\end{equation}
by comparison of the constant and \(q\) coefficients.
The product formula for \(\Delta\) shows that it has no zero in
the upper half-plane. The possible Coxeter numbers are
\[
 0,2,3,4,5,6,7,8,9,10,12,13,14,16,18,22,25,30,46.
\]
Their target values \(720-24h\) lie in \([-384,720]\).
If \(j(\tau)\) is rational at an imaginary quadratic argument,
the class polynomial has degree one. The complete list of these
rational singular moduli is given in
\cite[Table~1, p.~1057]{BLM2017}; the only member in this interval
is \(0\), of discriminant \(-3\).
Hence \(h=30\), and the Niemeier classification gives precisely
the two stated root systems. The equation \(j(\tau)=0\)
characterizes the modular orbit of \(\rho\).

On a block diagonal period matrix,
the definition of the genus-two theta series gives
\[
 \vartheta_2(L)(\operatorname{diag}(\tau_1,\tau_2))
       =\Theta_L(\tau_1)\Theta_L(\tau_2).
\]
For \(L=\Lambda\), neither factor is zero when both arguments
are imaginary quadratic. Every point of \(H_1\) has such a
block diagonal representative, and the Siegel modular transformation
factor is nonzero. The claimed nonvanishing follows.
This assertion concerns \(H_1\); it makes no claim away from that
surface or about the zeros of an Epstein zeta function.
\end{proof}

\section{Ogg's list and explicit class-number cutoffs}
Ogg's fixed-point formula already expresses the genus-zero question
in terms of quadratic class numbers. The additional assertion here is
an explicit finite cutoff, first elementary and then smaller under
the Dirichlet Riemann hypothesis.
Let \(h(D)\) denote the proper ideal class number of the quadratic
order of discriminant \(D<0\), and let \(\chi_D\) denote the
Kronecker character. For a prime \(p\ge5\), set
\[
 r(p)=\frac{h(-4p)+{\bf1}_{p\equiv3\ (4)}h(-p)}2,\qquad
 s(p)=\frac{p-1}{12}+\frac{1-\chi_{-4}(p)}4+
                         \frac{1-\chi_{-3}(p)}3.
\]

\begin{theorem}\label{thm:ogg-cutoffs}
For every prime \(p>7000\), one has \(r(p)<s(p)\), without any
Riemann hypothesis. Under the Dirichlet Riemann hypothesis the same
conclusion follows for every prime \(p>4300\).
The equality \(r(p)=s(p)\), for \(p\ge5\), holds exactly for
\[
 5,7,11,13,17,19,23,29,31,41,47,59,71.
\]
Together with \(2,3\), these are the primes for which
\(X_0(p)/\langle W_p\rangle\) has genus zero.
\end{theorem}
\begin{proof}
Ogg's Riemann--Hurwitz calculation
\cite[p.~452]{Ogg1974}, fixed-point formula
\cite[p.~454]{Ogg1974}, and genus formula
\cite[p.~455]{Ogg1974} give
\[
 g(X_0(p))=s(p)-1,\quad
 \#\operatorname{Fix}(W_p)=2r(p),\quad
 g(X_0(p)/\langle W_p\rangle)=\frac{s(p)-r(p)}2.
\]
Thus the class-number formulation itself is classical.

For a nonprincipal real primitive character of conductor \(q\),
periodicity and zero mean give
\[
 |L(1,\chi)|\le H_q+1\le\log q+2.
\]
Indeed truncate the Dirichlet series at \(q\), bound its initial
sum by \(H_q\), and integrate the tail by parts using
\(\left|\sum_{n\le x}\chi(n)\right|\le q\).
When \(p\equiv1\pmod4\), the class-number formula gives
\[
 r(p)=\frac{\sqrt p}{\pi}L(1,\chi_{-4p}).
\]
When \(p\equiv3\pmod4\), the order-conductor formula gives
\[
 h(-4p)=(2-\chi_{-p}(2))h(-p),\qquad
 r(p)\le \frac{2\sqrt p}{\pi}L(1,\chi_{-p}).
\]
Since \(\log(4p)+2\le2(\log p+2)\), both cases imply
\begin{equation}\label{eq:ogg-elementary}
 r(p)\le\frac{2\sqrt p}{\pi}(\log p+2).
\end{equation}
At \(p=7000\),
\[
 \frac{p-1}{\sqrt p}-\frac{24}{\pi}(\log p+2)>0.73815.
\]
The derivative of the left side is positive for \(p\ge7000\):
it is at least \(1/(2\sqrt p)-24/(\pi p)>0\).
Consequently \eqref{eq:ogg-elementary} is strictly smaller
than \((p-1)/12\le s(p)\) beyond that cutoff.

For the smaller cutoff, put
\[
 U(q)=2e^\gamma
 \left(\log\log q-\log2+\frac12+\frac1{\log\log q}\right).
\]
Under the Dirichlet Riemann hypothesis, the explicit estimate of
Lamzouri--Li--Soundararajan, extended to every \(q\ge3\) by
Languasco--Trudgian \cite[Theorem~1]{LT2022}, gives
\(L(1,\chi)\le U(q)\).
For \(p\ge4300\), \(U(4p)\le2U(p)\).
To see this, put \(x=\log\log p>1\).
The difference \(\log\log(4p)-x\) is at most \(\log2\);
the function \(x-\log2+1/2+1/x\) has derivative between zero
and one there, and its value is greater than \(\log2\).
It follows that \(r(p)\le2\sqrt p\,U(p)/\pi\).
At \(p=4300\),
\[
 \frac{p-1}{\sqrt p}-\frac{24}{\pi}U(p)>0.19855.
\]
Its derivative is at least
\[
 \frac1{2\sqrt p}-\frac{48e^\gamma}{\pi p\log p}>0
 \qquad(p\ge4300).
\]
This proves the second cutoff. Both numerical inequalities can
equivalently be checked with rational interval bounds for
\(\pi,e^\gamma,\log7000,\log4300\).

Finally, enumerate primitive reduced positive binary quadratic forms
\([a,b,c]\), with
\[
 b^2-4ac=D,\quad |b|\le a\le c,\quad
 b\ge0\ \text{if } |b|=a\text{ or }a=c.
\]
The reduction inequalities give \(a\le\sqrt{|D|/3}\), so this is
a finite exact computation. For each of the \(898\) primes
\(5\le p\le7000\), counting these forms for the two discriminants
in \(r(p)\) gives the displayed equality list and no instance of
\(r(p)>s(p)\). The workbench program \emph{Root exact checks}
records every integer count; a separate enumeration gives the
same list. This completes the finite verification.
\end{proof}

\section{Class-number terms in the \(M_{11}\) multiplicities}
The class-number estimates apply to a specified contribution to a virtual
multiplicity; they do not bound a cusp contribution or force the
multiplicity to be positive.
We use the \(2\)-optimal virtual graded modules of
Cheng--Duncan--Mertens \cite[Theorem~4.2.5 and Sect.~4.3]{CDM2025},
with their nontrivial multiplier assignment at orders four and eight.
Let \(H_N(D)\) be their generalized Hurwitz class number for
\(\Gamma_0(N)\), and \(H=H_1\). Their normalization gives
\[
 H(D)=h(D)\quad(D<-4\text{ fundamental}),\qquad
 H(-3)=\tfrac13,\quad H(-4)=\tfrac12.
\]
This normalization must be fixed before applying the analytic
class-number formula.

For \(N=1,2,3,6\), let \(C_N(D)\) be the coefficient of \(2R_N\),
where
\[
 R_N=\frac{12}{\varphi(N)}
       \sum_{M\mid N}\mu(N/M)
       \frac{M}{[\mathrm{SL}_2(\Z):\Gamma_0(M)]}\,H_M.
\]
Here the equality denotes the corresponding generating functions,
or their coefficient at a fixed \(D\).
Define the class-number contribution by
\[
 A(D)=\frac{C_1(D)}{144}-\frac{C_2(D)}{48}
              +\frac{C_3(D)}{18}-\frac{C_6(D)}6.
\]

\begin{theorem}\label{thm:m11-class}
Let \(D=f^2D_0<0\) be a discriminant, with \(D_0\) fundamental,
\(f\ge1\), and \(3\nmid D\). Put \(q_0=|D_0|\) and
\begin{equation}\label{eq:m11-B}
 B(D)=\chi_{D_0}(3)\bigl(1-\chi_{D_0}(2)\bigr)
 \prod_{\substack{\ell^e\parallel f\\\ell\ {\nm odd}}}
 \left(1+(\ell-\chi_{D_0}(\ell))
                  \frac{\ell^e-1}{\ell-1}\right).
\end{equation}
Then
\[
 A(D)=\frac{\sqrt{q_0}}{\pi}\,B(D)L(1,\chi_{D_0}).
\]
For \(q\ge3\), define
\begin{align*}
 U(q)&=2e^\gamma
 \left(\log\log q-\log2+\frac12+\frac1{\log\log q}\right),\\
 V(q)&=\frac{\pi^2}{12e^\gamma}
 \left(\log\log q-\log2+\frac12+\frac1{\log\log q}
                  +\frac{14\log\log q}{\log q}\right)^{-1}.
\end{align*}
Alp\"oge's Dirichlet Riemann hypothesis and the explicit
Lamzouri--Li--Soundararajan bounds imply
\begin{equation}\label{eq:m11-bounds}
 \frac{|B(D)|\sqrt{q_0}}{\pi}V(q_0)
 \le |A(D)|\le
 \frac{|B(D)|\sqrt{q_0}}{\pi}U(q_0).
\end{equation}
If \(B(D)\ne0\), \(A(D)\) has its sign; if \(B(D)=0\), \(A(D)=0\).
For fundamental \(D\), this reduces to
\[
 A(D)=\chi_D(3)(1-\chi_D(2))H(D).
\]
\end{theorem}
\begin{proof}
Substitution of \(N=1,2,3,6\) gives
\[
 A=-2H_1+H_2+2H_3-H_6.
\]
Separate binary quadratic forms of discriminant \(D\) according to
their content \(d\mid f\), writing \(g=f/d\) and \(D'=g^2D_0\).
For a primitive form of discriminant \(D'\), the number of isotropic
lines modulo a prime \(\ell\) is \(1+\chi_{D'}(\ell)\).
If \(\ell\mid d\), the form after multiplication by \(d\) is zero
modulo \(\ell\), so all \(\ell+1\) lines are allowed.
Thus the contribution of this content to \(H_N(D)\), for squarefree
\(N\), is the weighted order class number \(H_{\nm prim}(D')\)
times
\[
 \prod_{\ell\mid N}\nu_\ell,\qquad
 \nu_\ell=
 \begin{cases}\ell+1&\ell\mid d,\\
 1+\chi_{D'}(\ell)&\ell\nmid d.
 \end{cases}
\]
This follows by counting the cosets of \(\Gamma_0(N)\), whose
local quotients are \(\mathbb P^1(\F_\ell)\); stabilizer weights
give the same formula for the two exceptional unit groups.
Consequently the multiplier in \(A\) is
\[
 -2+\nu_2+2\nu_3-\nu_2\nu_3=(\nu_3-1)(2-\nu_2).
\]
Since \(3\nmid D\), \(\nu_3-1=\chi_{D_0}(3)\).
The order class-number formula, including unit factors, is
\[
 H_{\nm prim}(g^2D_0)=
 \frac{\sqrt{q_0}}{\pi}L(1,\chi_{D_0})\,
 g\prod_{\ell\mid g}\left(1-\frac{\chi_{D_0}(\ell)}{\ell}\right).
\]
For every odd \(\ell^e\parallel f\), summing the local terms gives
\[
 1+\sum_{j=1}^e\ell^j(1-\chi_{D_0}(\ell)/\ell)
 =1+(\ell-\chi_{D_0}(\ell))\frac{\ell^e-1}{\ell-1}.
\]
For the prime two put \(a=\chi_{D_0}(2)\).
If \(2\nmid f\), the multiplier is \(1-a\).
If \(2^e\parallel f\), \(e\ge1\), the terms with
\(v_2(g)<e\) have multiplier \(-1\), and the last term has
multiplier \(1\). Their sum is
\[
 2^{e-1}(2-a)-
 \left(1+(2-a)\sum_{j=1}^{e-1}2^{j-1}\right)=1-a.
\]
This proves \eqref{eq:m11-B} and the exact formula for \(A\).
The analytic class-number formula is positive, so it also proves
the assertion about the sign. Finally
\cite[Theorem~1.5]{LLS2015} gives the explicit upper and lower
bounds for \(L(1,\chi)\) at sufficiently large conductor;
\cite[Theorem~1]{LT2022} extends these same bounds to all
primitive conductors \(q\ge3\). Applying these bounds proves
\eqref{eq:m11-bounds}.
\end{proof}

To specify the remaining terms for even discriminants, define
\[
 T_4(D)=\sum_{a^2+b^2+c^2=|D|/4}(-1)^{b+c},\qquad
 T_8(D)=\sum_{a^2+4b^2+4c^2=|D|/4}(-1)^{b+c},
\]
where the sums are over integers, and are zero if \(4\nmid D\).
Let \(a_D=C_{8|4}(D)\) be the coefficient of the normalized
cuspidal Jacobi form \(\varphi_{8|4}\) of
\cite[equation~(4.2.9)]{CDM2025}.
Absorb its fixed contribution in the base trace function into the
parameter \(\kappa\): if \(\lambda\) denotes the coefficient
parameter in \cite[equation~(4.3.1)]{CDM2025}, put
\(\kappa=\Re(\lambda)-4\).
The parameter lattice is unchanged by this translation.

\begin{proposition}\label{prop:m11-full}
With this choice of parameter the full multiplicity is
\[
 m_{55}(W_D)=A(D)+\frac14T_4(D)-\frac12T_8(D)
                       +\frac{\kappa}{4}a_D.
\]
For odd fundamental \(D\), this becomes
\[
 m_{55}(W_D)=
 \chi_D(3)(1-\chi_D(2))H(D)+\frac{\kappa}{4}a_D.
\]
\end{proposition}
\begin{proof}
The character inner product in
\cite[equation~(4.3.1)]{CDM2025} gives
\[
 m_{55}=A-\frac18C_4+\frac14C_8+
                          \frac{\Re(\lambda)}4a_D.
\]
The explicit trace functions in
\cite[equation~(4.2.11)]{CDM2025} give
\(C_4=-2T_4\) and \(C_8=-2T_8-4a_D\).
For example,
\(\theta^0_{4,0}-\theta^0_{4,4}
=\sum_{b\in\Z}(-1)^bq^{b^2}\), while the corresponding
difference for index \(16\) is
\(\sum_b(-1)^bq^{4b^2}\).
Multiplication by \(\theta^0_{1,0}\theta_{1,0}(\tau,z)\)
then gives exactly the two ternary sums above.
Substitution proves the first formula. For odd \(D\) the two
ternary sums vanish, and Theorem~\ref{thm:m11-class} gives the second.
\end{proof}

In particular, a two-sided bound for \(A(D)\) is a bound for
the multiplicity after these explicitly specified terms are subtracted.
When \(D\) is odd and fundamental, its class-number multiplier is
\(0\), \(2\), or \(-2\); reversing the sign reverses the order of
the bounds for the signed quantity.

\section{Central values and the size of a multiplicity}
Tunnell's formula determines the square of the coefficient in this
case, including the dependence on the twist; it does not determine
its sign or give a distribution law for the family.
Put
\[
 \frac{\eta(4\tau)^5\eta(16\tau)}
      {\eta(2\tau)^2\eta(8\tau)}
       =\sum_{n\ge1}a(n)q^n,\qquad
 \beta=\int_1^\infty\frac{dx}{\sqrt{x^3-x}}.
\]

\subsection{The exact normalization of the Tunnell coefficient}
\label{sec:tunnell-normalization}
Write
\[
 \frac{\eta(4\tau)^5\eta(16\tau)}{\eta(2\tau)^2\eta(8\tau)}
   =\sum_{n\geq1}a(n)q^n,
 \qquad
 \varphi_{8|4}(\tau,z)=\sum_{n,r}C_{\varphi}(r^2-4n)q^ny^r,
\]
where $\varphi_{8|4}=(y^{-1}+y)q+O(q^2)$ is the normalization in
Cheng--Duncan--Mertens \cite[equation~(4.2.9)]{CDM2025}.
The exact comparison in this normalization is
\begin{equation}\label{eq:tunnell-factor-two}
 C_{\varphi}(-n)=\frac12a(n)\qquad(n\equiv3\pmod8).
\end{equation}
In particular $C_{\varphi}(-3)=1$ whereas $a(3)=2$.
The factor two does not affect a nonvanishing criterion, but does affect
an identity of coefficients and any explicit constant multiplying their
squares.

\begin{proof}
First we identify the Jacobi form itself:
\begin{equation}\label{eq:phi84-explicit}
 \varphi_{8|4}(\tau,z)
   =\eta(2\tau)\eta(8\tau)^2\theta_{1,1}(\tau,z).
\end{equation}
Here is a modular verification, in addition to the leading-term check.
The Jacobi triple product gives
$\theta^0_{1,1}=2\eta(4\tau)^2/\eta(2\tau)$, so the right side is
\[
 2\eta(4\tau)^2\eta(8\tau)^2
    \frac{\theta_{1,1}(\tau,z)}{\theta^0_{1,1}(\tau)}.
\]
The eta criterion puts $\eta(4\tau)^2\eta(8\tau)^2$ in
$S_2(\Gamma_0(32))$ with trivial character: its two congruence sums are
$2\cdot4+2\cdot8=24$ and $2\cdot(32/4)+2\cdot(32/8)=24$,
and $4^2 8^2$ is a square.  The theta transformation law shows that
$\theta_{1,1}$ transforms by a scalar on $\Gamma_0(4)$, with the same
scalar at $z=0$; their ratio has weight zero and index one there.
It is holomorphic in the upper half-plane, since $\theta^0_{1,1}$ has
no zeros there.  Cusp holomorphy can be checked without dividing:
the positive eta product $\eta(2\tau)\eta(8\tau)^2$ has positive order
at every cusp, and the transformed index-one theta vector has only
nonnegative theta exponents.  Their product is therefore a Jacobi cusp
form of weight two and index one for $\Gamma_0(32)$.
Its first term is $(y^{-1}+y)q$.  The uniqueness specified in
Cheng--Duncan--Mertens \cite[equation~(4.2.9)]{CDM2025} proves \eqref{eq:phi84-explicit}.

Put $\vartheta(\tau)=\sum_{r\in\mathbb Z}q^{r^2}$.
Another triple-product identity gives
\[
 \frac{\eta(4\tau)^5\eta(16\tau)}{\eta(2\tau)^2\eta(8\tau)}
   =\vartheta(2\tau)\eta(8\tau)\eta(16\tau).
\]
The eta product on the right has exponents congruent to one modulo eight.
The terms with exponent congruent to three modulo eight come precisely
from the odd integers in $\vartheta(2\tau)$.  As
\[
 \sum_{r\ \mathrm{odd}}q^{2r^2}
    =2\frac{\eta(32\tau)^2}{\eta(16\tau)},
\]
we obtain the exact formal-series identity
\begin{equation}\label{eq:tunnell-sieve}
 \sum_{n\equiv3\ (8)}a(n)q^n=2\eta(8\tau)\eta(32\tau)^2.
\end{equation}
On the other hand the odd theta coefficient of
\eqref{eq:phi84-explicit}, with $\tau$ replaced by $4\tau$, is
\[
 \sum_{n\equiv3\ (8)} C_{\varphi}(-n)q^n
   =\eta(8\tau)\eta(32\tau)^2.
\]
Comparison with \eqref{eq:tunnell-sieve} proves
\eqref{eq:tunnell-factor-two} for every such $n$, without a finite
coefficient extrapolation.
\end{proof}

The coefficient $a(n)$ here is exactly the one in Tunnell's original
definition $g\vartheta_2$, where
$g=q\prod_{r\geq1}(1-q^{8r})(1-q^{16r})$
\cite[p.~323]{Tunnell1983}.  Thus the correction is a factor between the
two normalizations, not a change to Tunnell's definition.


\begin{theorem}\label{thm:tunnell-multiplicity}
Let \(D<0\) be squarefree with \(D\equiv21\pmod{24}\),
write \(n=|D|\), and let \(E_n:y^2=x^3-n^2x\).
For the parameter \(\kappa\) defined above one has
\[
 m_{55}(W_D)=\frac{\kappa}{8}a(n),\qquad
 L(E_n,1)=\frac{\beta}{4\sqrt n}\,a(n)^2.
\]
Consequently
\[
 m_{55}(W_D)^2=
       \frac{\kappa^2\sqrt n}{16\beta}L(E_n,1),\qquad
 |m_{55}(W_D)|=
       \frac{|\kappa|}{4\sqrt\beta}\,
       n^{1/4}\sqrt{L(E_n,1)}.
\]
If \(\kappa\ne0\), the zero multiplicities correspond exactly to
zero central values. For any \(X,M>0\), their exact counting
relation is
\[
 \begin{split}
 &\#\{n\le X:\ n\ {\nm squarefree},\ n\equiv3\ (24),\
                    |m_{55}(W_{-n})|\le M\}\\
 &\quad=
 \#\{n\le X:\ n\ {\nm squarefree},\ n\equiv3\ (24),\
             \sqrt n\,L(E_n,1)\le16\beta M^2/\kappa^2\}.
 \end{split}
\]
\end{theorem}
\begin{proof}
Here the discriminant is odd and divisible by three, so the
class-number multiplier in the direct calculation of
Proposition~\ref{prop:m11-full} is zero, and both ternary sums
vanish. The coefficient identity is \(a_D=a(n)/2\), as proved in
\eqref{eq:tunnell-factor-two}.
Tunnell's Theorem~3 \cite[pp.~328--329]{Tunnell1983}
gives the displayed central-value formula with this eta quotient
and this real period. Squaring the coefficient formula and
substituting gives the two identities. The counting relation is
their pointwise equivalence. If \(\kappa=0\), all these
multiplicities vanish and the claimed equivalence is inapplicable.
\end{proof}

A Riemann hypothesis for \(L(E_n,s)\) would place its nontrivial
zeros on \(\Re s=1\), which includes the central point itself.
It would also imply the conductor-aspect Lindel\"of estimate
\[
 L(E_n,1)\ll_\varepsilon n^\varepsilon,\qquad
 |m_{55}(W_{-n})|\ll_{\kappa,\varepsilon}n^{1/4+\varepsilon}.
\]
Indeed, the analytic conductor of the unitary normalization
$L(E_n,s+1/2)$ is $O(n^2)$, and the GRH upper bound of
Chandee \cite[Corollaries~1.1 and~1.3, p.~3]{Chandee2009}
is $C^{o(1)}$ as the conductor $C$ tends to infinity.
Substitution in the central-value identity gives the second estimate.
This is an upper bound, not a distribution theorem.
It would therefore neither force these multiplicities to be nonzero
nor determine their signs or their frequencies of occurrence.
The displayed counting identity is the precise transfer of any
independently proved central-value distribution to their absolute
values. A statement about ranks or congruent numbers requires its
own input beyond the location of zeros.


\section{The positive-infinity-type Hecke question}
Let $K=\Q(\omega)$, $\omega^2+\omega+1=0$, and let $\psi_{36}$ and
$\psi_{27}$ be the algebraic Hecke characters of
$y^2=x^3+1$ and $y^2+y=x^3$. The question is whether every nontrivial
zero of $L(\psi_{36},s)$ and $L(\psi_{27},s)$ has real part $1$, and
whether every nontrivial zero of $L(\psi_{36}^3,s)$ has real part $2$.
Those three Riemann hypotheses are not proved in this paper.

On a $j=0$ curve over $K$, the automorphism
$(x,y)\mapsto(\omega x,y)$ is defined over $K$ and commutes with
Frobenius at every prime ideal of good reduction. At an inert rational
prime, the residue field over $K$ has order $p^2$, so the same assertion
holds. This settles the local order-three action. A separate note gives
a prime-ideal criterion obtained from Alp\"oge's abstract analytic
argument. Its arithmetic hypothesis, an estimate for weighted prime-ideal
sums, remains unproved. The finite colour of the character records its
infinity type modulo six and does not determine its archimedean type.

\section*{Note on method}
The mathematics and text were produced by GPT-6 Astra (ChatGPT Pro mode).
The accompanying records distinguish exact computations, targeted proof
reading, and re-proofs from the statement by separate instances of the
same model. Prior checks attributed to other models are not presented as
checks of this paper. No human has checked the text. The authors dedicate
their original text and programs to the public domain under CC0~1.0.
The workbench is \url{\workbenchurl}; the permanent record index is
\url{\recordurl}.

\bibliographystyle{amsplain}
\bibliography{references}
\end{document}
